% ============================================================
%  Honey In The Cloud — Final Year Project Report
%  BSc (Hons) Computer Science & Information Technology
%  University of Galway
% ============================================================

\documentclass[12pt, a4paper]{report}   % 'report' class suits multi-chapter docs

% ---------- Encoding & Fonts ----------
\usepackage[T1]{fontenc}
\usepackage[utf8]{inputenc}
\usepackage{lmodern}
\usepackage{footnote}
\usepackage{microtype}

% ---------- Page Layout ----------
\usepackage[
  top=2.5cm, bottom=2.5cm,
  left=3cm,  right=2.5cm,
  headheight=14pt
]{geometry}

% ---------- Headers & Footers ----------
\usepackage{fancyhdr}
\pagestyle{fancy}
\fancyhf{}
\fancyhead[L]{\small\leftmark}
\fancyhead[R]{\small Honey In The Cloud}
\fancyfoot[C]{\thepage}
\renewcommand{\headrulewidth}{0.4pt}

% ---------- Colour ----------
\usepackage[dvipsnames]{xcolor}
\definecolor{linkblue}{RGB}{0,70,140}

% ---------- Hyperlinks ----------
\usepackage[
  colorlinks=true,
  linkcolor=linkblue,
  urlcolor=linkblue,
  citecolor=linkblue,
  pdftitle={Honey In The Cloud},
  pdfauthor={Chukwuma George Anayo-Ezikeoha},
  pdfsubject={BSc Final Year Project},
  pdfkeywords={honeypot, CDN, WAF, Cloudflare, BunnyCDN, HDBSCAN, cybersecurity}
]{hyperref}

% ---------- Mathematics ----------
\usepackage{amsmath, amssymb}

% ---------- Tables ----------
\usepackage{booktabs}
\usepackage{longtable}
\usepackage{array}
\usepackage{multirow}
\usepackage{tabularx}

% ---------- Graphics ----------
\usepackage{graphicx}
\graphicspath{{vertopal_0e99b85c758c4238b17cd9f53834280c/media/}}
\usepackage{float}
\usepackage[font=small, labelfont=bf, labelsep=period]{caption}

% ---------- Lists ----------
\usepackage{enumitem}
\setlist{noitemsep, topsep=4pt}

% ---------- Code Listings ----------
\usepackage{listings}
\lstset{
  basicstyle=\ttfamily\small,
  breaklines=true,
  frame=single,
  backgroundcolor=\color{gray!8},
  keywordstyle=\color{blue},
  commentstyle=\color{green!50!black},
  stringstyle=\color{orange!80!black},
  numbers=left,
  numberstyle=\tiny\color{gray},
  numbersep=6pt
}

% ---------- Miscellaneous ----------
\usepackage{parskip}          % space between paragraphs instead of indent
\usepackage[normalem]{ulem}
\usepackage{setspace}
\onehalfspacing               % 1.5 line spacing (typical for reports)
\usepackage{pdfpages}
\usepackage{emptypage}        % suppress headers on blank pages

% ---------- Chapter / Section Style ----------
\usepackage{titlesec}
\titleformat{\chapter}[display]
  {\normalfont\Large\bfseries}{\chaptertitlename\ \thechapter}{12pt}{\huge}
\titlespacing*{\chapter}{0pt}{-10pt}{20pt}

\providecommand{\frontmatter}{\pagenumbering{roman}}
\providecommand{\mainmatter}{\pagenumbering{arabic}}
\providecommand{\backmatter}{}

% ============================================================
%  DOCUMENT
% ============================================================
\begin{document}

% ============================================================
%  TITLE PAGE
% ============================================================
\begin{titlepage}
  \centering
  \vspace*{1cm}

  \includegraphics[width=0.55\textwidth]{image1.png}\\[1.5cm]

  {\LARGE\bfseries Honey In The Cloud}\\[0.6cm]
  {\large Comparative CDN and WAF Security Analysis with Unsupervised Attack Classification\\
   and Agentic Incident Response using Cloud-Based Honeypots}\\[2cm]

  \begin{tabular}{rl}
    \textbf{Programme:}  & BSc (Hons) Computer Science \& Information Technology \\[4pt]
    \textbf{School:}     & School of Computer Science, University of Galway \\[4pt]
    \textbf{Document:}   & Final Year Project Report \\[4pt]
    \textbf{Date:}       & 1 April 2026 \\[4pt]
    \textbf{Author:}     & Chukwuma George Anayo-Ezikeoha \\[4pt]
    \textbf{Student ID:} & 22358526 \\[4pt]
    \textbf{Supervisor:} & Dr. Enda Barrett
  \end{tabular}

  \vfill
\end{titlepage}

% ============================================================
%  FRONT MATTER  (roman page numbers)
% ============================================================
\frontmatter

% ---------- Abstract ----------
\chapter*{Abstract}
\addcontentsline{toc}{chapter}{Abstract}

Modern web applications face growing exposure to automated bots, credential
stuffing, and volumetric attacks. Cloud providers market CDNs, WAFs, and bot
mitigation services as comprehensive defences, yet independent empirical
evaluation of these claims, particularly on non-enterprise service tiers
remains scarce.

This project deploys a cloud-hosted honeypot disguised as a legitimate streaming
service (\texttt{hdtvstreams.com}) to attract real-world malicious traffic and
evaluate two CDN providers under identical conditions: Cloudflare (free tier)
and BunnyCDN (Advanced Shield). Testing revealed that while both platforms
absorbed volumetric network-layer attacks effectively, HTTP/2
protocol-compliant floods bypassed Cloudflare's edge in a significant proportion
of cases. Rate limiting enforced thresholds but permitted indefinite attack
cycling, and reduced telemetry access on lower service tiers consistently
limited incident visibility.

To complement CDN testing, an unsupervised machine-learning pipeline using
HDBSCAN clustering was trained on honeypot logs to classify live traffic into
behavioural attack profiles; identifying patterns ranging from PHP webshell
deployment and path traversal probing to secrets harvesting, IoT exploitation,
and Exchange/RDP reconnaissance. An agentic LLM layer interprets these clusters,
augments its context through Retrieval-Augmented Generation, and auto-generates
structured incident tickets. Critically, the same layer propagates derived
threat intelligence back into the firewall ruleset, enabling the system to
harden its own defences in response to observed attack patterns growing more
sophisticated as the threats it encounters do.

The findings expose a measurable gap between vendor security promises and
real-world effectiveness at lower service tiers, and demonstrate that a
lightweight, closed-loop detection and response system can provide continuous,
self-reinforcing protection without enterprise-scale infrastructure.


% ---------- Acknowledgements ----------
\chapter*{Acknowledgements}
\addcontentsline{toc}{chapter}{Acknowledgements}

I would like to thank my project supervisor, Dr. \ Enda Barrett, for his guidance and support throughout
the duration of the project. A huge thank you to the School of Computer Science at
the University of Galway for providing the academic framework within which this
research could be conducted.

I am grateful to the open-source communities behind Locust, hping3, slowhttptest,
h2load, Scikit-learn and golden-eye, whose tools made the experimental work possible.
I extend my gratitude to Google Cloud, Amazon Web Services, Cloudflare and BunnyCDN
for providing accessible free and trial tiers that enabled independent security
evaluation without enterprise licensing.

Thank you to God for giving me the discipline, creativity, and fortitude to persevere
through blockers and complete this project.

% ---------- Table of Contents ----------
\tableofcontents
\listoffigures
\listoftables

% ============================================================
%  MAIN MATTER  (arabic page numbers)
% ============================================================
\mainmatter

% ============================================================
%  CHAPTER 1 — INTRODUCTION
% ============================================================
\chapter{Introduction}

\section{Project Description}

One way in which defenders can learn an attacker's strategy is to dupe them into
attacking a honeypot. The attacker believes they are attacking a production system,
but in actual fact it is a fake entity designed to trap the attacker. As the attacker
attempts to compromise the honeypot, we can learn about their attack strategy: what
vectors they are using, what commands they are executing and in what order. Learning
the attack motivation and strategy enables cybersecurity experts to harden their
defences~\cite{awesome-honeypots}. If an attacker realises they are interacting with 
a honeypot they will immediately stop; building a system that emulates a real system 
convincingly is an open challenge. The goal is to keep them interacting while also 
learning their strategies.

\section{Useful Definitions}
\label{sec:definitions}

\begin{description}[leftmargin=2cm, style=nextline]

  \item[Honeypot]
    A deliberately vulnerable system designed to attract and trap attackers. It mimics
    a legitimate production environment to deceive malicious actors, allowing defenders
    to observe attack strategies and techniques without risking real 
    systems~\cite{serverless-honeypot}.

  \item[Content Delivery Network (CDN)]
    A geographically distributed network of servers that cache and deliver web content
    from locations closer to end users. Modern CDNs also provide DDoS protection, bot
    mitigation, and web application firewalls~\cite{cf-how-works}.

  \item[Web Application Firewall (WAF)]
    A security solution that filters, monitors, and blocks HTTP/HTTPS traffic to and
    from web applications based on predefined security rules. WAFs protect against SQL
    injection, XSS, and other OWASP Top 10 vulnerabilities~\cite{owasp-waf}.

  \item[CAPTCHA]
    A challenge-response test designed to determine whether a user is human or an
    automated bot. Common implementations include image-selection puzzles, text
    recognition, or invisible background checks~\cite{cf-captcha-end}.

  \item[Bot]
    An automated software application that performs repetitive tasks over the internet.
    Bots can be benign (search engine crawlers) or malicious (credential stuffing
    tools, web scrapers, vulnerability scanners)~\cite{cf-verified-bots}.

  \item[DoS (Denial of Service)]
    An attack that attempts to make a machine, network, or service unavailable to its
    intended users by overwhelming it with illegitimate requests, exhausting resources
    such as CPU, memory, and bandwidth~\cite{stallings2003dos}.

  \item[DDoS (Distributed Denial of Service)]
    A cyberattack that attempts to overwhelm a target system, service, or network with
    excessive traffic from multiple sources, rendering it unavailable to legitimate
    users~\cite{cf-ddos}.

  \item[Adaptive Deception Technology]
    Security systems that dynamically adjust their behaviour and responses based on
    attacker interactions, making honeypots more realistic and harder to detect while
    gathering intelligence on attack strategies~\cite{awesome-honeypots}.

  \item[Serverless Computing]
    A cloud computing model where the cloud provider manages server infrastructure,
    automatically allocating resources as needed. Developers deploy code without
    managing underlying servers, paying only for actual compute time 
    consumed~\cite{aws-lambda}.

  \item[API (Application Programming Interface)]
    A set of protocols and tools that allows different software applications to
    communicate with each other. In web contexts, APIs enable front-end applications
    to interact with back-end services through defined 
    endpoints~\cite{aws-gateway}.

  \item[Rate Limiting]
    A security control that restricts the number of requests a user or IP address can
    make to a service within a specified time period, helping prevent abuse, scraping,
    and denial-of-service attacks~\cite{cf-rate-limiting}.

  \item[Telemetry]
    The automated collection and transmission of data from remote sources for
    monitoring and analysis. In security contexts, telemetry includes logs, metrics,
    and event data used to understand system behaviour and detect 
    threats~\cite{aws-cloudwatch}.

  \item[Challenge Page]
    An intermediate page presented to users before granting access to a website,
    typically containing a CAPTCHA or JavaScript verification to distinguish humans
    from automated bots~\cite{cf-bot-fight}.

  \item[Brute Force Attack]
    A hacking method that involves using trial and error to crack passwords or gain
    unauthorised access~\cite{fortinet-brute-force}.

\end{description}

\section{Introduction}

Modern web applications face increasingly sophisticated threats: automated bots,
web scrapers, credential-stuffing campaigns, and opportunistic actors constantly
probing for exploitable weaknesses~\cite{siteguarding}. Cloud providers have 
responded with integrated security offerings; Content Delivery Networks (CDNs) 
with built-in filtering, Web Application Firewalls (WAFs), and bot-mitigation 
services~\cite{cf-reverse-proxy}. Companies such as Cloudflare, Google Cloud, 
and Amazon Web Services market these as comprehensive defensive layers, capable 
of identifying and blocking malicious traffic before it reaches production 
systems~\cite{cf-how-works,aws-cloudfront}.

Whether these promises hold under real-world conditions is a harder question.
Independent evaluation of commercial CDN security features, especially on the
free and lower-paid tiers accessible to most developers, is limited~\cite{cloudflare-wp}. 
Do rate-limiting rules meaningfully constrain sustained attack traffic? Is it
feasible to build a system that learns from the attacks it observes and
autonomously tightens its own posture in response?

This project addresses these questions through empirical analysis using adaptive
deception technology. A cloud-based honeypot deployed on Google Cloud Platform,
fronted by alternating CDN configurations, was disguised as a legitimate
streaming service to attract genuine attack traffic creating a controlled but
authentic environment in which vendor security claims can be measured against
observed outcomes.

Beyond CDN evaluation, the diversity of collected traffic motivated an automated
threat-classification and response pipeline. An unsupervised machine-learning
layer clusters behavioural patterns in the logs~\cite{sklearn-hdbscan}, while 
an agentic LLM component interprets those clusters, generates incident 
reports~\cite{jira-api}, and propagates threat intelligence directly into the 
system's firewall rules~\cite{cf-api-lockdown}, allowing defences to adapt 
continuously as the system accumulates knowledge of attacker behaviour.

The full source code and deployment configuration for this project are publicly
available at \url{https://github.com/G-cae78/ProtectTheHoney}.

\section{Motivation}

The motivation for this research stems from several observations about the current
state of cloud security:

\begin{description}[leftmargin=2.5cm, style=nextline]

  \item[Market-Reality Gap]
    Cloud providers invest heavily in marketing their security features, but
    independent verification of these claims is scarce~\cite{cloudflare-wp}. 
    There are no clear SLAs specifying how effective these security offerings 
    should be.

  \item[The CAPTCHA Paradox]
    End users are increasingly burdened with CAPTCHA challenges, sometimes multiple
    times per session. It is not fully verified whether CAPTCHAs are security theatre
    that inconveniences legitimate users while barely impacting malicious 
    automation~\cite{cf-captcha-end}.

  \item[Configuration Complexity]
    Modern CDNs like Cloudflare offer dozens of toggleable security settings. However,
    it is unclear which combination actually enhances security versus which are
    redundant~\cite{cf-bot-fight}. This research aims to isolate and test individual 
    features.

  \item[Multi-Cloud Comparisons]
    Organisations often employ multiple cloud providers, yet there is limited
    comparative data on how their security offerings stack up against identical
    threats~\cite{azure-frontdoor,aws-cloudfront}. Does Cloudflare's bot detection 
    outperform AWS WAF? Is Google Cloud CDN's scraping protection more or less 
    effective?

  \item[Evolution of Attack Patterns]
    CDN attack-intelligence models may be trained on older patterns~\cite{concept-drift}; 
    this project captures contemporary attack patterns against realistic web 
    applications.

\end{description}

\section{Research Questions}
\label{sec:research-questions}

\begin{enumerate}

  \item \textbf{How effective are CDN-provided security mechanisms at detecting and
    mitigating real-world malicious traffic when deployed on non-enterprise tiers?}

  \item \textbf{How does the effectiveness of security features differ between
    Cloudflare, AWS CloudFront/WAF, and Google Cloud CDN when exposed to identical
    attack scenarios?}

  \item \textbf{To what extent do CAPTCHAs and JavaScript challenges reduce automated
    bot traffic without degrading legitimate user access?}

  \item \textbf{Which categories of malicious traffic (e.g.\ brute force, API abuse,
    automated scanning) are most likely to bypass CDN security controls?}

  \item \textbf{What is the rate of false positives (legitimate interactions blocked
    or challenged), and how does this differ across CDN platforms and settings?}

\end{enumerate}

Ultimately, this project aims to shift the conversation around cloud security from
trust-based adoption to evidence-based evaluation. By providing transparent,
reproducible measurements of what actually stops attacks, we enable security
professionals to make informed decisions rather than relying on vendor 
assurances~\cite{gartner-cf}.

% ---------- Related Work ----------
\chapter{Related Work}
\label{sec:related-work}

\section{Cloudflare CDN Performance: A WordPress Case Study~\cite{cloudflare-wp}}
\begin{figure}[H]
  \centering
  \includegraphics[width=0.8\textwidth]{image2.png}
  \caption{Cloudflare Network Architecture}
  \label{fig:cloudflare-arch}
\end{figure}
This research compared two WordPress websites with similar traffic volumes (25--33k
monthly visitors) hosted on the same provider, with one integrated with Cloudflare
CDN and the other using standard hosting. The Cloudflare-integrated site scored 95\%
on performance metrics compared to 51\% for the non-CDN site, with significant
improvements in time-to-first-byte, start render, and speed index. Memory usage also
differed dramatically: the Cloudflare site peaked at 400\,MB and remained stable,
while the non-CDN site repeatedly hit a 2\,GB memory limit. The study further noted
that Cloudflare's edge caching reduced the volume of requests reaching the origin
server by a significant margin, effectively lowering exposure to volumetric
traffic spikes. While the paper focuses on performance rather than security,
it establishes a useful empirical baseline for understanding how CDN interception
behaves under real-world traffic, a foundation this project builds upon by
evaluating Cloudflare's security controls against genuine malicious traffic rather
than legitimate user requests.

\section{Hybrid Deep Learning for Cloud Intrusion Detection~\cite{eos-ids}}
An Optimal Security System for intrusion detection in cloud computing environments
(EOS-IDS) was proposed using a three-stage process: (1) data pre-processing via an
Improved Heap Optimisation (IHO) algorithm; (2) feature selection via a Chaotic Red
Deer Optimisation (CRDO) algorithm; and (3) detection and classification via a Deep
Kronecker Neural Network (CKNN). The model successfully addressed common IDS
limitations including over-fitting, low classification accuracy, and high false
positive rates. Evaluation was conducted on benchmark datasets including NSL-KDD and
UNSW-NB15, where the system achieved notably high detection accuracy across multiple
attack categories such as DoS, probe, and R2L. Despite these results, the reliance
on pre-labelled, static datasets means the model has no mechanism for discovering
novel or unlabelled attack patterns, a limitation this project addresses through
the use of unsupervised clustering on live honeypot data.

\section{Intrusion Detection Using Hybrid Deep Learning~\cite{rana2024}}
This MSc research by Kunal Rana investigated ML and deep learning for intrusion
detection using the UNSW-NB15 dataset (nine attack types across 49 features). The
study combined traditional ML (logistic regression, decision trees) with LSTM,
bidirectional LSTM, and autoencoders. Bi-LSTM combined with autoencoders significantly
outperformed traditional methods, emphasising the importance of algorithm selection
and feature extraction strategies. The research also highlighted that the choice of
feature engineering pipeline had a substantial impact on model performance, with
raw feature sets performing considerably worse than those refined through systematic
selection. Notably, the study acknowledged that models trained on the UNSW-NB15
dataset may not generalise to traffic patterns observed in different deployment
environments, a shortcoming that motivates the live data acquisition strategy adopted
in this project.

\section{Concept Drift Detection via Pre-Clustering~\cite{concept-drift}}
This work on concept drift detection via pre-clustering and statistical testing
directly informs the drift monitoring approach used in the HDBSCAN pipeline
(Section~\ref{sec:drift}). The authors demonstrate that grouping data into clusters
prior to applying statistical drift tests substantially improves detection sensitivity
compared to applying tests on raw feature distributions. Their proposed framework
computes distributional divergence between reference and evaluation windows and
triggers retraining when drift exceeds a configurable threshold. This project
operationalises a closely related mechanism by monitoring centroid displacement
in UMAP projection space between successive evaluation windows, adapting the
underlying principle to the high-dimensional, sparse nature of HTTP request logs.

\medskip
Taken together, these works highlight two consistent gaps in the literature: the lack
of empirical CDN security evaluation under real attack conditions, and the predominance
of supervised or benchmark-independent approaches to traffic classification. This
project addresses both by combining live honeypot deployment with unsupervised
learning on naturally accumulated traffic data.

\section{Gaps in the Literature}
\label{sec:gaps}

Taken together, the surveyed works highlight several consistent gaps that this
project directly addresses.

\textbf{Absence of empirical CDN evaluation under real attack conditions.}
The majority of CDN and WAF research either relies on vendor-provided benchmarks or
evaluates synthetic attack scenarios in controlled laboratory environments. The
WordPress CDN comparison~\cite{cloudflare-wp} measures performance metrics rather
than security effectiveness, and no study was found that independently evaluates
free-tier CDN security claims against genuine malicious traffic attracted to a live
deployment. This project addresses the gap by deploying a real honeypot that passively
attracts organic attack traffic, then measuring what CDN controls actually intercept
versus what reaches the origin.

\textbf{Dominance of supervised learning on static benchmark datasets.}
Both IDS approaches reviewed~\cite{eos-ids,rana2024} train and evaluate on fixed
datasets (UNSW-NB15 and similar), which cannot reflect the continuously evolving
attack patterns seen on live infrastructure. Supervised models also require labelled
training data, which presupposes knowledge of attack categories before deployment.
This project uses HDBSCAN, an unsupervised algorithm, trained directly on live
honeypot logs with no pre-labelled ground truth, enabling it to discover novel attack
clusters without prior knowledge of their existence.

\textbf{Lack of closed-loop automated response.}
None of the surveyed works close the loop between threat detection and remediation.
Detection systems produce alerts or classifications, but the translation of those
outputs into actionable firewall rules and structured incident records remains a
manual analyst task. This project contributes an agentic LLM layer that autonomously
interprets cluster outputs, deploys Cloudflare WAF rules via the Rulesets API, and
raises structured Jira tickets, eliminating the analyst in the loop for routine
threat categories while decribing reasoning behind actions and grounding recommendations in the MITRE D3FEND
knowledgebase~\cite{d3fend} via RAG.

\textbf{No evaluation of concept drift in live security deployments.}
The concept drift work~\cite{concept-drift} develops theoretical detection mechanisms
but does not apply them to a live security pipeline. This project operationalises
drift monitoring by computing centroid displacement scores between training and test
UMAP projections after each model evaluation window, providing a practical trigger
for retraining without requiring manual inspection of cluster assignments.

%  CHAPTER 3 — SYSTEM DESIGN & ARCHITECTURE
% ============================================================
\chapter{System Design and Architecture}

\begin{figure}[H]
  \centering
  \includegraphics[width=\textwidth]{image3.png}
  \caption{Honeypot Service --- High-Level Architecture}
  \label{fig:arch}
\end{figure}

\section{Honeypot Design}
\label{sec:honeypot-design}

A legitimate domain (\url{https://hdtvstreams.com}) was purchased and equipped with a
valid SSL/TLS certificate. The certificate significantly increased traffic to the
honeypot: many automated scanning tools and botnets prioritise HTTPS-enabled sites
because they appear more legitimate~\cite{siteguarding} and are more likely to contain
valuable data. The presence of a real domain rather than a bare IP address also
makes the honeypot visible to broader DNS-based bot reconnaissance 
sweeps~\cite{mitre-recon}.

Within days of deploying the domain and certificate, traffic increased noticeably,
demonstrating how adding security measures (while essential for protecting legitimate
services) paradoxically makes a honeypot more effective by making it indistinguishable
from a real production application.

The honeypot is deployed on a Google Cloud Platform (GCP) virtual machine instance.
To ensure 24/7 availability without requiring an open SSH session, NGINX acts as a
reverse proxy and static file server~\cite{nginx-docs}, while PM2 (a production 
process manager for Node.js) keeps the SQLite backend running as a daemon, 
automatically restarting on crash.

The application also uses an AWS Lambda function to log interaction events through
the frontend to CloudWatch~\cite{aws-cloudwatch}. When triggered, the Lambda records 
basic request details (source IP, request headers, event type) and checks the IP 
against public reputation and blacklist sources~\cite{aws-lambda}.

\section{CDN-Fronted Architecture}

All incoming traffic is routed through the CDN before reaching the 
honeypot~\cite{cf-reverse-proxy}. This allows direct observation of how requests are 
handled at the CDN layer: blocked outright, challenged, or forwarded to the backend 
depending on the active configuration.

Suspicious requests that reach the application are logged using a serverless pipeline
triggered via an API gateway~\cite{aws-gateway}. Metadata including request paths, 
headers, IP addresses, user agents, and timestamps is recorded and stored centrally 
for later analysis~\cite{aws-cloudwatch}. This ensures consistent data collection 
across all test phases. The CDN remains the primary security control under evaluation, 
while the honeypot acts as a passive observation point.

\section{Software Tools}

\subsection{Development Tools}

\begin{itemize}
  \item Visual Studio Code
  \item BurpSuite~\cite{burpsuite}
  \item Postman
\end{itemize}

\subsection{Amazon Web Services (AWS)}

\begin{description}[leftmargin=3cm, style=nextline]

  \item[Lambda]
    Serverless computing service that executes logging code on trigger, records attack
    attempts, checks IP addresses against blacklists, and captures request
    metadata~\cite{aws-lambda}.

  \item[API Gateway]
    Creates the HTTP endpoint embedded in the Vue.js frontend and acts as the trigger
    for the Lambda function. The hidden endpoint captures bot interactions that normal
    users would not trigger~\cite{aws-gateway}.

  \item[CloudWatch]
    Stores all Lambda execution logs, which are reviewed to explain malicious traffic
    patterns and other honeypot interactions~\cite{aws-cloudwatch}.

  \item[EC2]
    Hosts the honeypot web application. Also served as the initial deployment target
    before migration to GCP~\cite{aws-ec2}.

  \item[CloudFront / WAF]
    AWS's CDN and Web Application Firewall. Used as a third CDN for comparison
    against Cloudflare and Google Cloud CDN~\cite{aws-cloudfront,azure-waf-rules}.

\end{description}

\subsection{Google Cloud Services}

\begin{description}[leftmargin=3cm, style=nextline]

  \item[Cloud Domains]
    Used to manage DNS records for the \texttt{hdtvstreams.com} domain, pointing
    traffic to the honeypot infrastructure.

  \item[Cloud CDN]
    Google's globally distributed CDN, tested as a further comparison point against
    Cloudflare's security features~\cite{cf-how-works}.

  \item[Compute Engine (GCP VM)]
    Hosts the production honeypot; replaced the original EC2 instance in January 2026.

\end{description}

\subsection{Cloudflare}

Primary CDN under evaluation, routing all traffic through their network to
empirically test the effectiveness of bot detection, CAPTCHA challenges, DDoS
protection, and other marketed security features against real-world 
attacks~\cite{cf-how-works,cf-ddos,cf-bot-fight}.

\subsection{Attack \& Penetration Testing Tools}

All offensive tests were run from a Kali Linux virtual machine:

\begin{itemize}
  \item \textbf{hping3} --- network packet manipulation tool~\cite{kali-hping3}
  \item \textbf{slowhttptest} --- application-layer slow HTTP attack 
        simulator~\cite{kali-slowhttp}
  \item \textbf{h2load} --- HTTP/2 benchmarking and load testing 
        tool~\cite{nghttp2-h2load}
  \item \textbf{GoldenEye} --- HTTP DoS testing tool~\cite{kali-goldeneye}
  \item \textbf{SQLMap} --- automated SQL injection testing tool~\cite{owasp-waf}
  \item \textbf{Ghauri} --- WAF-bypass-focused SQL injection framework~\cite{owasp-waf}
  \item \textbf{Locust} --- open-source Python load testing framework~\cite{locust}
\end{itemize}

\section{Preliminary Research Findings}

\subsection{Traffic Observations}

\begin{itemize}
  \item The hidden API endpoint was mostly hit by automated scrapers and LLM bots ---
        surprisingly non-malicious in nature~\cite{cf-verified-bots}.
  \item Real attacks targeted backend endpoints directly, skipping the frontend
        entirely~\cite{mitre-recon}.
  \item Mass exploit scans probing for exposed Laravel/PHPUnit endpoints were
        detected early~\cite{exploit-db}.
  \item IP reputation checks worked well to flag suspicious sources~\cite{aws-lambda}.
\end{itemize}

\subsection{Cloudflare Free Tier --- Baseline}

\begin{itemize}
  \item Gathered baseline data on what actually gets blocked versus what passes
        through~\cite{cf-how-works}.
  \item Tracked CAPTCHA frequency and whether bots could bypass 
        challenges~\cite{cf-captcha-end}.
  \item Recorded when rate limits triggered~\cite{cf-rate-limiting}.
  \item Tested whether different security settings produced measurably different
        results~\cite{cf-bot-fight}.
\end{itemize}

\subsection{Attack Patterns Identified Early}

\begin{itemize}
  \item Automated scanners probing for vulnerable frameworks~\cite{exploit-db}.
  \item Login brute-force attempts~\cite{fortinet-brute-force}.
  \item Random bot crawling with varied user agents~\cite{cf-verified-bots}.
  \item Direct API attacks that completely ignored the 
        frontend~\cite{owasp-api8}.
\end{itemize}

% ============================================================
%  CHAPTER 4 — RESEARCH & EXPERIMENTATION
% ============================================================
\chapter{Research and Experimentation}
\label{chap:experimentation}

\section{Methodology}
\label{sec:methodology}

This project follows an experimental, empirical research approach by observing how
real-world attacks interact with CDN security controls. Rather than relying solely on
documentation or synthetic attack tools, the method focuses on measuring how CDN
features behave when exposed to genuine malicious traffic from the public internet.

\begin{figure}[H]
  \centering
  \begin{minipage}{0.48\textwidth}
    \centering
    \includegraphics[width=\textwidth]{image7.png}
  \end{minipage}\hfill
  \begin{minipage}{0.48\textwidth}
    \centering
    \includegraphics[width=\textwidth]{image8.png}
  \end{minipage}
  \caption{Cloudflare Service Plans and Offerings}
  \label{fig:cloudflare-plans}
\end{figure}

Cloudflare was selected as the primary CDN platform due to its popularity and the
availability of a well-featured free tier~\cite{cf-how-works}; many small and 
medium-sized organisations rely on this tier in production, making it a realistic 
evaluation target. The scope is therefore limited to security features available 
without enterprise licensing~\cite{gartner-cf}. The following components are tested:

\begin{itemize}
  \item DDoS mitigation
  \item Web Application Firewall (WAF) with managed rules
  \item Bot protection mechanisms, including JavaScript challenges and CAPTCHA
  \item Basic rate limiting and request filtering
\end{itemize}

\section{Cloudflare}

\subsection{Setup and Configuration}

Setting up Cloudflare to front the honeypot was straightforward. The process began
by adding the domain to Cloudflare, which automatically scanned and displayed existing
DNS records~\cite{cf-how-works}. After verifying that these records pointed to the 
GCP instance IP address (\texttt{34.51.232.105}), ownership verification was completed 
by updating the domain's nameservers at the registrar (GCP Cloud Domains / Squarespace) 
to point to Cloudflare's nameservers.

\begin{figure}[H]
  \centering
  \begin{minipage}{0.48\textwidth}
    \centering
    \includegraphics[width=\textwidth]{image9.png}
  \end{minipage}\hfill
  \begin{minipage}{0.48\textwidth}
    \centering
    \includegraphics[width=\textwidth]{image10.png}
  \end{minipage}
  \caption{Default Name Servers vs Cloudflare Name Servers}
  \label{fig:nameservers}
\end{figure}

Once the nameserver changes propagated, Cloudflare began proxying traffic to the
honeypot. DNS lookups now return Cloudflare's edge IPs
(\texttt{172.67.213.165} and \texttt{104.21.23.209}) rather than the origin GCP IP.

\subsection{Rate Limiting with Locust}
\label{sec:cf-rate-limiting}

\subsubsection{Overview}

To validate Cloudflare's rate-limiting protections, load testing was conducted using
Locust~\cite{locust}, an open-source Python-based performance testing framework chosen
for its developer-friendly, code-based test definitions and gevent-based concurrency,
enabling a single process to simulate thousands of concurrent users.

\subsubsection{Test Setup}

A single simulated user repeatedly attempted to log into the honeypot's authentication
endpoint (\texttt{/api/auth/login}) using randomly generated credentials from the
Faker library~\cite{faker}. The wait time between requests was set to 10--30\,ms to
simulate an aggressive automated attack~\cite{fortinet-brute-force}. HTTP~429 
responses were treated as successful outcomes (rate limiter functioning correctly), 
and the script captured and logged the attacking IP (\texttt{140.203.12.65}) and the 
exact request count before each trigger.

\subsubsection{Results}

The Cloudflare rate-limiting rule was configured to trigger after 100 requests from a
single IP within a 10-second window, blocking that IP for 10 seconds (the maximum
available on free tier)~\cite{cf-rate-limiting}. The first trigger occurred at 
approximately 108 requests, validating the threshold.

\subsubsection{Security Implications}
\label{sec:cf-rate-implications}

While the initial 100-request burst was blocked, the same IP could immediately
resume attack cycles after each 10-second timeout. Over 24 hours, this theoretically
permits up to 864,000 credential attempts (100 requests every 10 seconds = 36,000 per
hour). Cloudflare does not natively support permanent or escalating blocks based on
repeated violations; custom solutions using the Cloudflare API are required for
persistent blocking~\cite{cf-community-blocking}.

\subsubsection{Potential Solutions}

\begin{enumerate}
  \item Increase the maximum block duration (up to 1 hour via WAF rules interface) to
        reduce daily attempts from hundreds of thousands to at most 2,400.
  \item Implement progressive penalties: trigger a 1-hour block after five consecutive
        violations of the 10-second window rule.
  \item Build a custom solution using Cloudflare's API~\cite{cf-api-lockdown} in 
        conjunction with a Lambda function~\cite{aws-lambda} to permanently block IPs 
        that repeatedly exceed the threshold.
\end{enumerate}

\subsection{DDoS Protection}
\label{sec:cf-ddos}

\subsubsection{Overview}

Cloudflare holds a 4.5-star rating based on 246 verified reviews on Gartner Peer
Insights~\cite{gartner-cf}. Its DDoS protection operates across multiple layers; the
two primary categories relevant to this testing are:

\begin{itemize}
  \item \textbf{SSL/TLS DDoS protection} --- application-layer attacks targeting SSL
        handshakes and HTTP/HTTPS services.
  \item \textbf{Network-layer DDoS protection} --- Layer 3/4 volumetric attacks (TCP 
        SYN floods, UDP floods, ICMP floods). Mitigated automatically at Cloudflare's 
        edge using its 477\,Tbps network capacity across 330+ 
        cities~\cite{cf-ddos}.
\end{itemize}

\begin{figure}[H]
  \centering
  \includegraphics[width=\textwidth]{image18.png}
  \caption{Cloudflare DDoS Protection Tiers}
  \label{fig:cf-ddos-tiers}
\end{figure}

\subsubsection{TCP SYN Flood (Layer 4)}

\paragraph{Overview}

A TCP SYN flood exploits the TCP three-way handshake by sending a high volume of SYN
packets with spoofed source IPs, resulting in half-open connections. The server waits
for final ACKs that never arrive, exhausting its connection table and preventing
legitimate connections~\cite{cf-syn-flood,tcp-handshake}.

\begin{figure}[H]
  \centering
  \includegraphics[width=0.5\textwidth]{image20.png}
  \caption{TCP SYN Flood Exploitation}
  \label{fig:syn-exploit}
\end{figure}

\paragraph{Attack Setup}

\begin{itemize}
  \item \textbf{Target:} Honeypot on GCP, fronted by Cloudflare.
  \item \textbf{Tool:} hping3 on Kali Linux VM~\cite{kali-hping3}.
  \item \textbf{Monitoring:} Nginx server logs~\cite{nginx-docs}, Cloudflare 
        real-time analytics, GCP observability dashboard.
\end{itemize}

\begin{lstlisting}[language=bash, caption={hping3 TCP SYN Flood Commands}]
sudo hping3 -S --flood -p 80 --rand-source <target-IP>
sudo hping3 -S -i u100 -p 80 hdtvstreams.com
\end{lstlisting}

\paragraph{Results}

\begin{figure}[H]
  \centering
  \includegraphics[width=0.85\textwidth]{image22.png}
  \caption{Resource Utilisation During SYN Flood}
  \label{fig:syn-resources}
\end{figure}

CPU utilisation remained unaffected; the regular spikes visible in the dashboard are
due to pre-existing scripts. Network traffic showed a minor increase, peaking at only
17.21\,KiB/s at the origin, negligible for a production-grade service. SYN-ACK
responses were initially returned from Cloudflare's edge IP (\texttt{172.67.213.165}),
confirming the handshake was handled at the edge, not the origin. As the flood
intensified, intermittent ``ICMP Network Unreachable'' messages appeared, suggesting
additional filtering was triggered.

\paragraph{Analysis}

Cloudflare absorbed and mitigated the SYN flood before it could affect the origin
server~\cite{cf-syn-flood}. The stable CPU and negligible network utilisation at the 
origin confirm effective edge mitigation. However, the lack of telemetry on the free 
tier makes it difficult to diagnose root causes or understand the exact mitigation path.

\subsubsection{UDP Flood (Layer 4)}

\paragraph{Overview}

UDP is a connectionless ``fire-and-forget'' protocol: unlike TCP, no acknowledgement
packets are involved. A UDP flood sends massive packet volumes without maintaining
connection state, targeting random ports and forcing the server to generate ICMP
``Destination Unreachable'' responses for each, consuming bandwidth and 
CPU~\cite{cf-udp-flood,udp-wiki}.

\begin{figure}[H]
  \centering
  \begin{minipage}{0.48\textwidth}
    \centering
    \includegraphics[width=\textwidth]{image26.png}
  \end{minipage}\hfill
  \begin{minipage}{0.48\textwidth}
    \centering
    \includegraphics[width=\textwidth]{image27.png}
  \end{minipage}
  \caption{UDP Spoofing vs Normal UDP Interaction}
  \label{fig:udp-comparison}
\end{figure}

\paragraph{Attack Setup}

\begin{itemize}
  \item \textbf{Target:} Honeypot on GCP, fronted by Cloudflare.
  \item \textbf{Tool:} hping3 on Kali Linux VM~\cite{kali-hping3}.
\end{itemize}

\begin{lstlisting}[language=bash, caption={hping3 UDP Flood Commands}]
sudo hping3 -2 --flood -p 80 --rand-source <target-IP>
sudo hping3 -2 -i u100 --spoof -p 80 hdtvstreams.com
sudo hping3 -2 -i u100 -p 80 hdtvstreams.com
\end{lstlisting}

\paragraph{Results}

\begin{figure}[H]
  \centering
  \includegraphics[width=0.7\textwidth]{image28.png}
  \caption{Resource Utilisation During UDP Flood}
  \label{fig:udp-resources}
\end{figure}

CPU utilisation was completely unaffected. Network traffic showed only a minor
increase. Across all test variations, 100\% packet loss was observed --- not a single
UDP packet received a response indicating Cloudflare is silently dropping
non-DNS UDP traffic at the edge~\cite{cf-udp-flood}.

Notably, during one run the attacking machine was itself rate-limited for
approximately 30 seconds. Replacing the \texttt{--flood} flag with the \texttt{-i}
interval flag prevented further rate limiting while still permitting up to 1 million
UDP packets at 1k packets/second.

\paragraph{Analysis}

Silent edge-level drops mean network-layer UDP floods never generate dashboard events
on free-tier plans~\cite{cf-ddos}. While the protection functioned correctly, the 
complete absence of logging prevents post-incident analysis.

\subsubsection{SSL/TLS DDoS --- SlowHTTPTest (Layer 7)}

\paragraph{Overview}

Slow HTTP attacks exploit the way web servers handle concurrent connections: legitimate
HTTP connections are opened but data is sent intentionally at low rates, holding
connections open and exhausting the server's connection pool~\cite{kali-slowhttp}. 
SlowHTTPTest supports Slowloris (slow headers), Slow HTTP POST (slow body), and 
slow-read attack modes.

\paragraph{Attack Setup}

\begin{itemize}
  \item \textbf{Target:} Honeypot on GCP, fronted by Cloudflare.
  \item \textbf{Tool:} slowhttptest on Kali Linux VM~\cite{kali-slowhttp}.
\end{itemize}

\begin{lstlisting}[language=bash, caption={SlowHTTPTest Command}]
slowhttptest -c 1000 -H -g -o slow_reprt -i 10 r -r 200 -t GET 
  -u https://hdtvstreams.com -x 24 -p 3
\end{lstlisting}

\paragraph{Results}

\begin{figure}[H]
  \centering
  \includegraphics[width=0.65\textwidth]{image32.png}
  \caption{Resource Utilisation During Slow HTTP Attack}
  \label{fig:slowhttp-resources}
\end{figure}

CPU and network traffic were completely unaffected. With 1,000 connections attempted,
all connections were initially kept open for approximately 10 seconds before Cloudflare
began automatically closing them. When scaled to 5,000 connections, the defence
mechanism activated earlier.

\begin{figure}[H]
  \centering
  \begin{minipage}{0.48\textwidth}
    \centering
    \includegraphics[width=\textwidth]{image33.png}
  \end{minipage}\hfill
  \begin{minipage}{0.48\textwidth}
    \centering
    \includegraphics[width=\textwidth]{image34.png}
  \end{minipage}
  \caption{SlowHTTPTest Attack Progression --- 1,000 Connections}
  \label{fig:slowhttp-1000}
\end{figure}

\begin{figure}[H]
  \centering
  \begin{minipage}{0.48\textwidth}
    \centering
    \includegraphics[width=\textwidth]{image35.png}
  \end{minipage}\hfill
  \begin{minipage}{0.48\textwidth}
    \centering
    \includegraphics[width=\textwidth]{image36.png}
  \end{minipage}
  \caption{SlowHTTPTest Attack Progression --- 5,000 Connections}
  \label{fig:slowhttp-5000}
\end{figure}

Table~\ref{tab:slowhttp-phases} summarises the five mitigation phases observed:

\begin{table}[H]
  \centering
  \caption{SlowHTTPTest Attack Phases and Cloudflare Response (5,000 connections)}
  \label{tab:slowhttp-phases}
  \begin{tabularx}{\textwidth}{lXXXXX}
    \toprule
    \textbf{Seconds} & \textbf{Closed} & \textbf{Pending} & \textbf{Connected}
      & \textbf{Available} & \textbf{Phase} \\
    \midrule
    0  & 0     & 1 & 0     & 5000 & Baseline \\
    20 & 365   & 4 & 966   & 5000 & Rapid Connection Growth \\
    40 & 1731  & 4 & 1248  & 5000 & Sustained High Pressure \\
    60 & 3390  & 5 & 1296  & 5000 & Peak Load (Pre-Mitigation) \\
    \textbf{64} & \textbf{3730} & \textbf{0} & \textbf{1270} & \textbf{5000}
      & \textbf{Defence Trigger} \\
    80 & {\raise.17ex\hbox{$\scriptstyle\sim$}}4941 & 0 & 0 & 5000 
      & Attack Neutralised \\
    \bottomrule
  \end{tabularx}
\end{table}

\paragraph{Analysis}

Despite Cloudflare clearly responding with progressive connection termination and
eventual complete blocking, \textbf{no DDoS analytics events appeared in the
dashboard}~\cite{cf-ddos}. This silent-protection gap limits post-incident analysis 
and means administrators have no way to know when their service was under attack.

\subsubsection{HTTP/2 Flood --- H2Load (Layer 7)}
\label{sec:h2load-cf}

\paragraph{Overview}

H2Load is a benchmarking tool for HTTP/2, HTTP/1.1, and HTTP/3 protocols, developed
as part of the nghttp2 project~\cite{nghttp2-h2load}. Critically, it uses the
system's OpenSSL library, producing valid TLS fingerprints and legitimate SNI
extensions~\cite{http2-ddos}. This forces CDN systems to process requests rather than 
drop them as malformed packets, making H2Load a more realistic Layer 7 defence test.

\paragraph{Test Command}\mbox{}\\[4pt]

\begin{lstlisting}[language=bash, caption={H2Load Test Command}]
h2load -n 10000 -c 100 -m 10 https://hdtvstreams.com
\end{lstlisting}

\paragraph{Results}

\begin{figure}[H]
  \centering
  \includegraphics[width=0.4\textwidth]{image38.png}
  \caption{Resource Utilisation During H2Load}
  \label{fig:h2load-resources}
\end{figure}

CPU utilisation increased 3$\times$ and network traffic 500$\times$ during attack 
periods. The initial test (10,000 requests, 100 concurrent clients) achieved a 100\% 
success rate. When scaled to 100,000 requests, Cloudflare's rate limiting activated 
but only blocked 68,356 requests, 31,644 (31.6\%) still reached the origin. On a 
repeated identical run, 56,801 requests succeeded, demonstrating inconsistency in 
Cloudflare's mitigation response~\cite{http2-ddos}.

\begin{figure}[H]
  \centering
  \includegraphics[width=\textwidth]{image39.png}
  \caption{H2Load Request Statistics Across Test Runs}
  \label{fig:h2load-stats}
\end{figure}

\paragraph{Analysis}

Unlike previous attacks, H2Load generated visible security events in Cloudflare's
dashboard and triggered an automated notification email. However, the 31--60\% bypass rate across runs 
demonstrates that legitimate-protocol Layer 7 traffic poses a greater challenge to 
CDN-based DDoS protection than volumetric network-layer attacks~\cite{http2-ddos}.

\subsubsection{DDoS Attack Pattern Summary --- Cloudflare}
\label{sec:cf-ddos-summary}

\begin{table}[H]
  \centering
  \caption{DDoS Test Results Summary for Cloudflare Free Tier}
  \label{tab:cf-ddos-summary}
  \begin{tabularx}{\textwidth}{lllXXX}
    \toprule
    \textbf{Attack Type} & \textbf{Tool} & \textbf{Layer}
      & \textbf{Origin Impact} & \textbf{Edge Mitigation} 
      & \textbf{Disruption} \\
    \midrule
    SYN Flood    & hping3       & L4 & None       & Yes            & No \\
    UDP Flood    & hping3       & L4 & None       & Yes            & No \\
    ICMP Flood   & hping3       & L3 & None       & Yes            & No \\
    HTTP Flood   & GoldenEye    & L7 & Minimal    & Yes            & No \\
    Slow HTTP    & slowhttptest & L7 & Minimal    & Partial        & No \\
    SSL Flood    & thc-ssl-dos  & L7 & Low        & Partial        & No \\
    HTTP/2 Flood & h2load       & L7 & Noticeable & Less effective 
      & Latency / CPU spike \\
    \bottomrule
  \end{tabularx}
\end{table}

\subsection{SQL Injection Testing --- Cloudflare Free Tier}
\label{sec:cf-sqli}

\subsubsection{Overview}

SQL injection remains one of the most prevalent and dangerous web application
vulnerabilities, consistently appearing in the OWASP Top 10~\cite{owasp-waf}.
This test evaluated whether Cloudflare's free-tier WAF managed rules would detect
and block automated SQL injection tooling against the honeypot's login endpoint
(\texttt{/login} and \texttt{/api/auth/login}), using both the industry-standard
SQLMap scanner and the WAF-bypass-focused Ghauri framework. All tests were conducted
from a Kali Linux virtual machine.

\begin{itemize}
  \item \textbf{Attack OS:} Kali Linux
  \item \textbf{Tools:} SQLMap (v1.10.2) and Ghauri (v1.4.3)
  \item \textbf{Interception:} Burp Suite Community Edition~\cite{burpsuite}
  \item \textbf{Target:} \url{https://hdtvstreams.com} via Cloudflare Free Tier Edge
\end{itemize}

\subsubsection{Test 1 --- Automated Form Discovery}

\begin{lstlisting}[language=bash, caption={SQLMap Form Discovery}]
sqlmap -u https://hdtvstreams.com/login --forms --batch
\end{lstlisting}

\paragraph{Result}

SQLMap returned HTTP~405 (Method Not Allowed) responses throughout the scan. Without
a \texttt{--data} flag, the tool had no POST body to target and defaulted to probing
the \texttt{Referer} and \texttt{Host} headers instead of the intended form fields.
No injection point was identified. More significantly, \textbf{no Cloudflare WAF
event was triggered} and no block page was returned, the 405 response originated
entirely from the origin server's Nginx configuration~\cite{nginx-docs}, not from any 
CDN-layer filtering. The SQLMap default user-agent passed through the Cloudflare edge 
without challenge.

\subsubsection{Test 2 --- Aggressive Payload Injection}

\begin{lstlisting}[language=bash, caption={SQLMap Aggressive Test --- Maximum Risk and Level}]
sqlmap -u https://hdtvstreams.com/login \
  --data="email=test@example.com&password=test" \
  --risk 5 --level 5 \
  --batch
\end{lstlisting}

\paragraph{Result}

With \texttt{--data} supplied, SQLMap correctly targeted the \texttt{email} and
\texttt{password} POST body parameters. The \texttt{--level 5} flag additionally
enabled header injection testing across \texttt{User-Agent}, \texttt{Referer}, and
\texttt{Host}. SQLMap reported both POST parameters as ``not dynamic'', the server
returned an identical response body regardless of payload content depriving the
tool of the differential signal required for detection. The scan concluded with no
injectable parameters found.

Critically, \textbf{zero Cloudflare WAF security events were generated} across the
entire run. Risk~5 / Level~5 is the most exhaustive SQLMap configuration, producing
a high volume of highly characteristic SQL injection payloads in the POST body;
none were flagged, blocked, or challenged by Cloudflare's managed 
ruleset~\cite{owasp-waf}.

\subsubsection{Test 3 --- Evasion with Tamper Scripts and Random User Agent}

\begin{lstlisting}[language=bash, caption={SQLMap with Tamper Scripts and Random Browser UA}]
sqlmap -u https://hdtvstreams.com/login \
  --data="email=test@example.com&password=test" \
  --tamper=randomcase,space2comment,charencode \
  --random-agent \
  --batch
\end{lstlisting}

\paragraph{Result}

This test introduced two evasion techniques simultaneously: tamper scripts to
obfuscate payloads (case randomisation, comment-space substitution, and URL character
encoding), and \texttt{--random-agent} to replace SQLMap's identifiable default
user-agent with a randomised real browser string~\cite{cf-verified-bots}.

The \texttt{--random-agent} flag produced the most significant finding of the entire
Cloudflare SQL injection test series. With a browser-like user agent, all requests
bypassed the Cloudflare edge without any challenge and hit the origin server directly.
The Cloudflare traffic dashboard recorded \textbf{59,950 requests originating from a
single Irish IP address} (ASN~1213 --- HEAnet, \texttt{140.203.12.65}) with
\textbf{zero mitigation actions applied}. Every request detail entry showed
``Suspicious activity: None'' and ``Mitigation: Not mitigated''
(Figures~\ref{fig:cf-sqli-dashboard} and~\ref{fig:cf-sqli-detail}).

\begin{figure}[H]
  \centering
  \begin{minipage}{0.48\textwidth}
    \centering
    \includegraphics[width=\textwidth]{cf_sqli_source_ips.png}
    \caption*{Source IPs: 59.95k requests, single IP}
  \end{minipage}\hfill
  \begin{minipage}{0.48\textwidth}
    \centering
    \includegraphics[width=\textwidth]{cf_sqli_countries.png}
    \caption*{Country breakdown: Ireland 59.95k, France 1.01k, India 113}
  \end{minipage}
  \caption{Cloudflare Traffic Dashboard During SQLMap \texttt{--random-agent} Run}
  \label{fig:cf-sqli-dashboard}
\end{figure}

This result demonstrates that Cloudflare's free-tier WAF has no payload-level SQL
injection detection independent of user-agent heuristics~\cite{owasp-waf}. Switching 
from SQLMap's default UA to any browser string was sufficient to remove all friction.

\subsubsection{Test 4 --- Ghauri WAF-Bypass Framework}

\begin{lstlisting}[language=bash, caption={Ghauri with POST Request File and Random UA}]
ghauri -r post.txt --random-agent --batch --ignore-code 401
\end{lstlisting}

\paragraph{Result}

Ghauri initially returned 401 Unauthorized responses; adding \texttt{--ignore-code 401}
instructed the tool to proceed regardless. Ghauri tested a comprehensive range of
injection techniques --- error-based, boolean-based blind, time-based blind, and
stacked queries --- across both \texttt{email} and \texttt{password} fields. The run
concluded with:

\begin{lstlisting}
[CRITICAL] all tested parameters do not appear to be injectable.
\end{lstlisting}

As with all SQLMap runs, \textbf{no Cloudflare WAF rule fired} and no security event
appeared in the dashboard. Ghauri is specifically designed to bypass WAF 
fingerprinting~\cite{azure-waf-rules}; its evasion techniques proved unnecessary given 
that Cloudflare's free tier generated no detectable response to any of the injected 
payloads.

\subsubsection{Analysis}

Across all four test runs --- spanning basic form discovery, maximum-aggression payload
injection (\texttt{--risk 5 --level 5}), multi-technique tamper-script evasion, and a
dedicated WAF-bypass framework. Cloudflare's free-tier WAF generated
\textbf{zero security events} and applied \textbf{zero mitigations}.

Two independent factors explain why no successful injection occurred despite the
complete absence of WAF intervention:

\begin{enumerate}
  \item \textbf{Parameterised queries.} The Node.js backend uses SQLite with
        parameterised statements (\texttt{db.get("SELECT * FROM users WHERE email = ?
        AND password = ?", [email, password])}), which are structurally immune to SQL
        injection regardless of WAF coverage~\cite{nist-800-53}.
  \item \textbf{Uniform server responses.} The login endpoint returned an identical
        response body for all inputs, preventing injection tools from obtaining the
        differential signal required for blind injection detection.
\end{enumerate}

The critical security implication is that an identical application using 
string-concatenated queries would have been \textbf{fully exposed}: Cloudflare's 
free-tier managed rules provided no payload-level detection across any of the four 
scenarios~\cite{owasp-waf}. A simple \texttt{--random-agent} flag was sufficient to 
ensure 59,950 requests reached the origin with no challenge applied. This represents 
a meaningful gap in the free-tier WAF's out-of-the-box SQL injection coverage, with 
the application's own parameterised queries acting as the sole effective defence.

\begin{table}[H]
  \centering
  \caption{SQL Injection Test Results Summary --- Cloudflare Free Tier}
  \label{tab:cf-sqli-summary}
  \small
  \begin{tabularx}{\textwidth}{lXXX}
    \toprule
    \textbf{Test} & \textbf{Tool / Flags} & \textbf{WAF Response} 
      & \textbf{Outcome} \\
    \midrule
    Form discovery
      & SQLMap \texttt{--forms --batch}
      & No WAF events; 405 from Nginx
      & No injection found \\
    \addlinespace
    Aggressive injection
      & SQLMap \texttt{--risk 5 --level 5}
      & Zero WAF events; no block
      & Parameters non-injectable; uniform responses \\
    \addlinespace
    Tamper + random UA
      & SQLMap \texttt{--tamper --random-agent}
      & 59.95k requests unmitigated
      & All requests reached origin; zero challenges \\
    \addlinespace
    WAF-bypass tool
      & Ghauri \texttt{--random-agent --ignore-code 401}
      & No WAF events; no mitigation
      & Parameters non-injectable; parameterised 
        queries prevented exploitation \\
    \bottomrule
  \end{tabularx}
\end{table}

\subsection{WAF Rule Test --- Geographic Blocking}
\label{sec:cf-geo}

\subsubsection{Setup}

A custom WAF rule was implemented to block all requests originating from Ireland
(country code \texttt{IE})~\cite{cf-api-lockdown}.

\subsubsection{Results}

Accessing the honeypot from an Irish IP immediately returned Cloudflare's block page
in both standard and incognito modes. Programmatic testing via Locust~\cite{locust} 
confirmed HTTP 403 responses with Cloudflare-specific headers (\texttt{Server: 
cloudflare}, \texttt{CF-RAY: 9c1e5134afa8bde8-DUB}). The \texttt{-DUB} suffix 
indicates the request was handled at Cloudflare's Dublin PoP. Zero Nginx log entries 
were generated for blocked requests.

\begin{figure}[H]
  \centering
  \includegraphics[width=\textwidth]{image43.png}
  \caption{WAF Geographic Block Page Shown in Browser}
  \label{fig:waf-geo-block}
\end{figure}

\subsubsection{VPN Bypass}

Connecting through a VPN server outside Ireland immediately allowed access, with
requests appearing to originate from the VPN server's location~\cite{xvpn}.

\subsubsection{Analysis}

Geographic blocking effectively reduces attack surface from regions with no legitimate
users and is resource-efficient~\cite{cf-how-works}. However, IP-based location is a 
weak identifier for a determined adversary: VPN and proxy services trivially bypass 
it~\cite{xvpn}.

\section{BunnyCDN}
\label{sec:bunnycdn}

\subsection{Overview}

BunnyCDN (now Bunny.net) was founded in 2015 and is headquartered in Ljubljana,
Slovenia, giving it European jurisdiction relevant to GDPR and data 
sovereignty~\cite{european-tech-bunny}. Bunny.net operates 119 PoPs across six 
continents with capacity exceeding 150\,Tbps and an average global latency of 
25\,ms~\cite{european-tech-bunny}. Its security offering, Bunny Shield, combines WAF, 
DDoS mitigation, bot detection, rate limiting, access control lists, and upload 
scanning~\cite{bunny-shield-docs}.

For this project, the Advanced Shield tier was used (with \$20 free credit), providing
access to 255 base WAF rules, an AI-powered WAF learning mode, bot mitigation, curated
threat access lists, and upload scanning~\cite{bunny-shield-docs}.

\begin{figure}[H]
  \centering
  \includegraphics[width=0.7\textwidth]{image49.png}
  \caption{Bunny.net Subscription Tiers}
  \label{fig:bunny-tiers}
\end{figure}

\subsection{Setup and Configuration}

Setting up Bunny.net followed a process similar to Cloudflare: the domain was added
to Bunny's pull zone dashboard, nameservers were updated to \texttt{coco.bunny.net}
and \texttt{kiki.bunny.net}, and once propagated, DNS lookups returned Bunny.net's
edge IP (\texttt{87.249.137.52})~\cite{bunny-shield-docs}.

\begin{figure}[H]
  \centering
  \includegraphics[width=\textwidth]{image51.png}
  \caption{Updated Nameservers After Fronting with Bunny.net}
  \label{fig:bunny-ns}
\end{figure}

\subsection{Rate Limiting}
\label{sec:bunny-rate}

\subsubsection{Test Setup}

Bunny.net rate limiting was evaluated in three phases using Locust~\cite{locust} (500 
concurrent users, \texttt{/api/auth/login} POST requests and homepage GET requests). 
Tests were run both before and after a rule was configured to block for 30 seconds 
when an IP exceeded 300 requests per minute~\cite{bunny-shield-docs}. To isolate 
rate-limiting behaviour from bot fingerprinting (Locust's default Python user-agent is 
immediately challenged), a Python script using the \texttt{curl-cffi} library with the 
\texttt{chrome110} impersonation flag was also used.

\subsubsection{Results}

\begin{figure}[H]
  \centering
  \includegraphics[width=0.9\textwidth]{image53.png}
  \caption{Locust Request Results --- Before Rate Limiting Rule}
  \label{fig:bunny-before-rl}
\end{figure}

\begin{figure}[H]
  \centering
  \includegraphics[width=\textwidth]{image54.png}
  \caption{Locust Request Results --- After Rate Limiting Rule}
  \label{fig:bunny-after-rl}
\end{figure}

Without a rule, RPS climbed steadily to \textasciitilde380 and plateaued, with
\textasciitilde50\% failures from bot detection. After the 300 req/min rule was
applied, RPS built normally then collapsed sharply from \textasciitilde300 to near
zero almost immediately. P95 response times spiked to over 18,000\,ms, the 
signature of edge-level blocking rather than throttling~\cite{bunny-waf-docs}.

\begin{figure}[H]
  \centering
  \includegraphics[width=0.9\textwidth]{image55.png}
  \caption{Locust Chart --- Before Rate Limiting Rule}
  \label{fig:bunny-chart-before}
\end{figure}

\begin{figure}[H]
  \centering
  \includegraphics[width=\textwidth]{image56.png}
  \caption{Locust Chart --- After Rate Limiting Rule}
  \label{fig:bunny-chart-after}
\end{figure}

\begin{figure}[H]
  \centering
  \includegraphics[width=0.9\textwidth]{image57.png}
  \caption{Bunny Shield Dashboard Event Logs}
  \label{fig:bunny-events}
\end{figure}

\subsubsection{Analysis}

Unlike Cloudflare's 10-second maximum block on the free tier~\cite{cf-rate-limiting}, 
Bunny.net Advanced offers block durations from 30 seconds to 1 
hour~\cite{bunny-shield-docs}, giving administrators meaningful control over persistent 
attackers without custom automation. Additionally, Bunny Shield's managed rules 
classified randomly generated credentials containing SQL-like characters as MSSQL 
injection attempts, demonstrating more aggressive out-of-the-box detection than 
Cloudflare's free tier~\cite{bunny-waf-docs}.

\begin{figure}[H]
  \centering
  \includegraphics[width=0.85\textwidth]{image61.png}
  \caption{Bunny Shield Dashboard --- WAF Rule Violation Event}
  \label{fig:bunny-waf-event}
\end{figure}

\subsection{DDoS Protection}
\label{sec:bunny-ddos}

BunnyCDN's WAF provides five configurable DDoS sensitivity tiers~\cite{bunny-shield-docs}; 
testing was conducted at the ``Lowest'' level to establish baseline behaviour.

\begin{figure}[H]
  \centering
  \includegraphics[width=\textwidth]{image62.png}
  \caption{Bunny Shield DDoS Sensitivity Tiers}
  \label{fig:bunny-ddos-tiers}
\end{figure}

\subsubsection{TCP SYN Flood and UDP Flood (Layer 4)}

Both network-layer flood attacks were conducted using hping3~\cite{kali-hping3} with 
the same commands used against Cloudflare, allowing a direct side-by-side comparison.

\paragraph{TCP SYN Flood}
A SYN flood exploits the TCP three-way handshake by sending a high volume of SYN
packets, exhausting the target's connection table with half-open connections that
never complete~\cite{cf-syn-flood,tcp-handshake}.

\begin{lstlisting}[language=bash, caption={hping3 TCP SYN Flood Commands}]
sudo hping3 -S --flood -p 80 --rand-source <target-IP>
sudo hping3 -S -i u100 -p 80 hdtvstreams.com
\end{lstlisting}

\paragraph{UDP Flood}
UDP is connectionless and stateless; a flood sends massive packet volumes to random
ports, forcing the server to generate ICMP ``Destination Unreachable'' responses for
each, consuming bandwidth and CPU~\cite{cf-udp-flood,udp-wiki}.

\begin{lstlisting}[language=bash, caption={hping3 UDP Flood Commands}]
sudo hping3 -2 --flood -p 80 --rand-source <target-IP>
sudo hping3 -2 -i u100 --spoof -p 80 hdtvstreams.com
sudo hping3 -2 -i u100 -p 80 hdtvstreams.com
\end{lstlisting}

\paragraph{Results}

\begin{figure}[H]
  \centering
  \includegraphics[width=0.85\textwidth]{image63.png}
  \caption{hping3 Logs During SYN and UDP Flood Testing --- 100\% Packet Loss}
  \label{fig:bunny-hping-logs}
\end{figure}

Both attacks produced identical results across every run:

\begin{itemize}
  \item \textbf{100\% packet loss} --- not a single SYN or UDP packet received a
        response from either the edge or the origin.
  \item \textbf{Zero CPU and network impact} at the origin --- the GCP instance
        remained completely unaffected throughout both tests.
  \item At points during both runs, hping3 was \textbf{unable to resolve the target
        host entirely}, suggesting BunnyCDN applied temporary IP-level blocking on top
        of packet filtering --- a more aggressive response than silent dropping 
        alone~\cite{bunny-shield-docs}.
\end{itemize}

\paragraph{Analysis}

BunnyCDN absorbed both Layer 4 floods entirely at the edge, with no traffic reaching
the origin server. Unlike Cloudflare, which returned SYN-ACK responses before
beginning to filter~\cite{cf-syn-flood}, BunnyCDN dropped all packets silently from 
the first request indicating filtering occurs upstream of the TCP 
handshake~\cite{bunny-shield-docs}. The temporary host resolution failures observed 
during both floods suggest an additional IP-level throttling mechanism beyond standard 
packet filtering. However, as with Cloudflare's handling of network-layer attacks, 
\textbf{neither the SYN flood nor the UDP flood generated any dashboard events}, 
silent mitigation at this layer means administrators would have no visibility into the 
attack without external monitoring.

\subsubsection{Slow HTTP (SlowHTTPTest)}

The slowhttptest tool~\cite{kali-slowhttp} was executed in Slow Headers mode against 
the BunnyCDN-protected domain. The tool attempted to establish 1,000, 10,000, 200, 
and 65,535 concurrent slow HTTP connections, each sending partial headers over a 
240-second test window.

\paragraph{Resource Utilisation}

No requests reached the origin service. CPU utilisation at the origin was zero
attack traffic was absorbed entirely at the CDN edge before reaching the 
application~\cite{bunny-shield-docs}. Network traffic analysis confirmed the same: no 
slow-header packets reached the origin's network interface, indicating upstream 
filtering by BunnyCDN.

\paragraph{Attack Tool Observations}

The slowhttptest tool saves results to an HTML report and a CSV file, visualising how
connections behaved over time across each run. A clear pattern emerges across all
three runs: at approximately the 55--60 second mark, BunnyCDN begins aggressively
closing connections, and this behaviour repeats consistently regardless of the
connection count (1,000, 10,000, or 65,535). Before that threshold a small number of
connections hang in a connected state, but once the window closes they are cut off
rapidly. This suggests BunnyCDN enforces a hard timeout policy at the edge: connections
that have been sitting idle or sending partial headers for roughly a minute are
terminated, rather than being dropped immediately upon detection~\cite{bunny-waf-docs}.

\begin{figure}[H]
  \centering
  \includegraphics[width=0.7\textwidth]{image71.png}
  \caption{SlowHTTPTest Report --- Run 1}
  \label{fig:bunny-slow-run1}
\end{figure}

\begin{figure}[H]
  \centering
  \includegraphics[width=0.7\textwidth]{image72.png}
  \caption{SlowHTTPTest Report --- Run 2}
  \label{fig:bunny-slow-run2}
\end{figure}

\paragraph{Limitations}

Whilst testing BunnyCDN's connection limits, an attempt was made to open 1,000,000
concurrent connections; however the tool's built-in ceiling of 65,539 constrained the
maximum theoretical attack volume achievable in this environment. The true resilience
ceiling of BunnyCDN against larger-scale slow-header floods therefore remains untested.
The test was also conducted from a single source IP, which makes detection and blocking
significantly easier for a CDN than a distributed attack originating from many sources
would~\cite{stallings2003dos}.

\begin{figure}[H]
  \centering
  \includegraphics[width=0.85\textwidth]{image74.png}
  \caption{Console Error When Attempting More Than 65,539 Connections}
  \label{fig:bunny-slow-limit}
\end{figure}

\paragraph{BunnyCDN Protection}

Bunny Shield demonstrated effective mitigation of the slowhttptest attack. The
dashboard recorded 293 bot-detection events during the test window, most of which
were challenged, with zero requests attributed to DDoS 
mitigation~\cite{bunny-shield-docs}. This indicates that the traffic was classified 
as automated bot activity rather than a volumetric DDoS, and was handled accordingly. 
The combination of edge-level connection termination and bot-detection classification 
prevented any attack traffic from reaching the origin, resulting in no measurable 
impact on CPU or network resources.

\subsubsection{H2Load (HTTP/2 Flood)}
\label{sec:bunny-h2load}

The same H2Load configuration used against Cloudflare 
(Section~\ref{sec:h2load-cf}) was applied here~\cite{nghttp2-h2load}:

\begin{itemize}
  \item Requests: 10,000 / 100,000
  \item Concurrent Clients: 100
  \item Max Concurrent Streams: 10
  \item Protocol: HTTP/2 (TLS\_AES\_256\_GCM\_SHA384)
\end{itemize}

\paragraph{Test 1 --- Initial Probe}

2,800 of 10,000 requests succeeded, with 7,007 failed. TLS negotiation succeeded and
HTTP/2 was confirmed. Some filtering was already active.

\paragraph{Test 2 --- Escalated Detection}

Completed in 20.3\,s; only 912 of 10,000 requests succeeded. Connection times
averaged 238\,ms, suggesting throttling or challenge mechanisms~\cite{bunny-waf-docs}.

\paragraph{Test 3 --- Full Block}

100,000 requests: 0 succeeded, all returned 0 bytes. Connections were refused
immediately at network/firewall level before any HTTP response was generated.

\paragraph{Test 4 --- IP-level Rate Limiting}

Source IP was fully banned at infrastructure level; all connections dropped
instantly~\cite{bunny-shield-docs}.

\begin{figure}[H]
  \centering
  \includegraphics[width=0.85\textwidth]{image78.png}
  \caption{H2Load Console Errors After IP Block}
  \label{fig:bunny-h2load-errors}
\end{figure}

\begin{figure}[H]
  \centering
  \includegraphics[width=0.4\textwidth]{image79.png}
  \caption{GCP Instance Resource Utilisation During H2Load}
  \label{fig:bunny-h2load-resources}
\end{figure}

\begin{figure}[H]
  \centering
  \includegraphics[width=\textwidth]{image80.png}
  \caption{Bunny CDN Dashboard Response to H2Load Tests}
  \label{fig:bunny-h2load-dashboard}
\end{figure}

\subsubsection{Analysis}

Bunny.net Shield uses a progressive mitigation strategy: initial light filtering,
challenge mechanisms, connection throttling, and hard IP block~\cite{bunny-shield-docs}. 
This layered approach proved more effective than Cloudflare against H2Load, ultimately 
reaching a complete block by the fourth run~\cite{http2-ddos}.

\subsection{SQL Injection Testing}
\label{sec:bunny-sqli}

\subsubsection{Setup}

\begin{itemize}
  \item Attack OS: Kali Linux
  \item Tools: SQLMap (v1.10.2) and Ghauri (v1.4.3)
  \item Interception: Burp Suite Community Edition~\cite{burpsuite}
  \item Target: \url{https://hdtvstreams.com} via Bunny.net Edge
\end{itemize}

\subsubsection{Results}

\paragraph{Test 1 --- Automated Form Discovery (SQLMap)}\mbox{}\\[4pt]

\begin{lstlisting}[language=bash, caption={SQLMap Form Discovery}]
sqlmap -u https://hdtvstreams.com/login --forms --batch
\end{lstlisting}

SQLMap's user-agent signature was immediately identified via WAF Rule ID 913100:
``Found User-Agent associated with security scanner''~\cite{bunny-waf-docs}.

\paragraph{Test 2 --- Aggressive Payload Injection}\mbox{}\\[4pt]

\begin{lstlisting}[language=bash, caption={SQLMap Aggressive Test}]
sqlmap -u https://hdtvstreams.com/login \
  --data="email=test@example.com&password=test" \
  --risk 5 --level 5
\end{lstlisting}

WAF Rule ID 942230 fired: ``conditional SQL injection 
attempts''~\cite{azure-waf-rules}. JSON payloads containing nested CASE statements 
were caught and blocked.

\begin{figure}[H]
  \centering
  \includegraphics[width=\textwidth]{image85.png}
  \caption{WAF Log --- Captured SQLMap Payload}
  \label{fig:sqlmap-waf}
\end{figure}

\paragraph{Test 3 --- Evasion Attempts (Tamper Scripts)}\mbox{}\\[4pt]

\begin{lstlisting}[language=bash, caption={SQLMap with Tamper Scripts}]
sqlmap -u [URL] --data="[DATA]" \
  --tamper=randomcase,space2comment,charencode --batch
\end{lstlisting}

The libjection module (WAF Rule ID 942100) normalised the obfuscated traffic and
identified the underlying SQL injection pattern~\cite{bunny-waf-docs}.

\paragraph{Test 4 --- Ghauri with Burp Suite POST Request}\mbox{}\\[4pt]

\begin{lstlisting}[language=bash, caption={Ghauri Test Using Captured POST Request}]
ghauri -r post.txt --batch
\end{lstlisting}

No tool managed to identify the database type, let alone enumerate 
tables~\cite{owasp-waf}.

\begin{figure}[H]
  \centering
  \includegraphics[width=\textwidth]{image87.png}
  \caption{Request Interception with Burp Suite}
  \label{fig:burpsuite-intercept}
\end{figure}

\begin{figure}[H]
  \centering
  \includegraphics[width=\textwidth]{image88.png}
  \caption{Dashboard Showing Specific WAF Rule Infringements}
  \label{fig:bunny-waf-rules}
\end{figure}

\subsubsection{Analysis}

Bunny Shield demonstrated consistent, layered detection across all four test
scenarios~\cite{bunny-waf-docs}. In Log mode, all attacks were flagged; switching to 
Block mode terminated all flagged requests at the edge with no traffic forwarded to 
the origin.

\section{Security Comparison: Cloudflare vs BunnyCDN}
\label{sec:comparison}

\begin{table}[H]
  \centering
  \caption{Cloudflare Free Tier vs Bunny.net Advanced Shield --- Comparative 
  Results~\cite{cf-ddos,bunny-shield-docs,owasp-waf,http2-ddos}}
  \label{tab:comparison}
  \footnotesize
  \setlength{\tabcolsep}{4pt}
  \begin{tabularx}{\textwidth}{@{}l l X X X@{}}
    \toprule
    \textbf{Category} & \textbf{Attack} & \textbf{Cloudflare Free} 
      & \textbf{Bunny Advanced} & \textbf{Key Observations} \\
    \midrule

    Rate Limiting & Credential Stuffing
      & Triggers at 100 req/10\,s; max 10\,s block; indefinite 
        cycling~\cite{cf-rate-limiting}
      & Triggers at 300 req/min; 30\,s--1\,h blocks; hard 
        cut-off~\cite{bunny-shield-docs}
      & Bunny offers longer blocks and harder shutdown. \\
    \addlinespace

    L3 DDoS & ICMP Flood
      & Edge mitigation; no dashboard events~\cite{cf-ddos}
      & Edge mitigation; no detailed logs~\cite{bunny-shield-docs}
      & Both absorbed attack before origin impact. \\
    \addlinespace

    L4 DDoS & SYN Flood
      & Edge mitigation; 17\,KiB/s at origin; no 
        events~\cite{cf-syn-flood}
      & 100\% packet loss; temporary IP block; zero origin 
        impact~\cite{bunny-shield-docs}
      & Bunny more aggressive (source-level blocking). \\
    \addlinespace

    L4 DDoS & UDP Flood
      & Silent drop; 100\% loss; no events~\cite{cf-udp-flood}
      & 100\% loss; temporary source 
        blocking~\cite{bunny-shield-docs}
      & Both strong at network-layer filtering. \\
    \addlinespace

    L7 DDoS & Slow HTTP
      & Progressive mitigation; defence trigger 
        \textasciitilde64\,s; no dashboard alert~\cite{cf-ddos}
      & Edge timeout policy; 293 bot events 
        logged~\cite{bunny-waf-docs}
      & Cloudflare silent; Bunny classified as bot activity. \\
    \addlinespace

    L7 DDoS & HTTP/2 Flood
      & 31--60\% traffic reached origin; email + dashboard 
        alert~\cite{http2-ddos}
      & Progressive to full block by run 4; IP-level 
        ban~\cite{bunny-shield-docs}
      & Bunny escalated to complete block; Cloudflare 
        inconsistent. \\
    \addlinespace

    Bot Detection & Scripted Requests
      & Many passed without challenge~\cite{cf-bot-fight}
      & Immediate fingerprint detection; managed 
        challenge~\cite{bunny-waf-docs}
      & Bunny Shield more aggressive. \\
    \addlinespace

    WAF & SQL Injection
      & Zero WAF events across all 4 tests; \texttt{--random-agent} 
        bypassed all heuristics (59.95k 
        unmitigated)~\cite{owasp-waf}
      & Rule 913100 (scanner UA); Rule 942230 (payloads); 
        libjection normalised tamper obfuscation; all 4 tests 
        logged~\cite{bunny-waf-docs,azure-waf-rules}
      & Cloudflare free tier generated no SQLi events; Bunny 
        caught every test. \\
    \addlinespace

    Geo Blocking & Country Block
      & Edge-level 403; never reached origin; VPN bypass 
        trivial~\cite{xvpn}
      & Not tested
      & IP-based filtering effective but easily bypassed. \\
    \addlinespace

    Telemetry & ---
      & Limited logging; silent mitigation; email for HTTP DDoS 
        only~\cite{cf-ddos}
      & Granular WAF logs; clear rule violation 
        logging~\cite{bunny-waf-docs}
      & Bunny Advanced provides stronger forensic visibility. \\

    \bottomrule
  \end{tabularx}
\end{table}

% ============================================================
%  CHAPTER 5 — TRAFFIC PATTERN ANALYSIS
% ============================================================
\chapter{Traffic Pattern Analysis}
\label{chap:traffic}

\section{Overview}

The aim of this phase is to analyse live honeypot traffic in five-minute windows,
using these request buckets to group similar request patterns and generate actionable
security insights. By clustering logs, high-volume bot behaviour and targeted
malicious activity can be identified to prevent exploitation of potential
vulnerabilities~\cite{ml-clustering-sec}. An unsupervised clustering approach works 
in the right ``adaptive'' direction: it can highlight new anomalies as they arise and 
provide references to similar previous attacks, speeding up attack identification and 
surfacing new malicious behaviour~\cite{concept-drift}.

\section{Clustering Approach}

\subsection{Algorithm 1: K-Means Clustering}

\subsubsection{Overview}

K-Means is a popular unsupervised machine learning algorithm that partitions a dataset
into a pre-specified $k$ clusters, each defined by a centroid (the average position of
all its members)~\cite{sklearn-kmeans}. The algorithm iterates: randomly initialise 
$k$ centroids, assign each point to the nearest centroid, recalculate centroids, and 
repeat until convergence.

\subsubsection{Pre-processing}

Client IPs and timestamps were extracted from raw honeypot logs, then traffic was
aggregated into 5-minute time windows per IP. Features included request counts,
unique path counts, status code distributions (401, 403, 404, 200), and TF-IDF
vectors of request sequences. Numerical features were standardised using scikit-learn's
\texttt{StandardScaler}~\cite{sklearn-scaler}; TF-IDF vectors were horizontally
stacked. The dataset was split 70:30 (431 training / 185 test samples).

\subsubsection{Model Training}

The elbow method was used to determine optimal $k$ (tested 2--10 
clusters)~\cite{sklearn-kmeans}. The optimal value was identified as $k=5$.

\begin{figure}[H]
  \centering
  \begin{minipage}{0.45\textwidth}
    \centering
    \includegraphics[width=\textwidth]{image90.png}
  \end{minipage}\hfill
  \begin{minipage}{0.50\textwidth}
    \centering
    \includegraphics[width=\textwidth]{image91.png}
  \end{minipage}
  \caption{Elbow Method and K-Means Training Result Summary}
  \label{fig:kmeans-elbow}
\end{figure}

\begin{figure}[H]
  \centering
  \includegraphics[width=0.7\textwidth]{image92.png}
  \caption{K-Means PCA Cluster Visualisation}
  \label{fig:kmeans-pca}
\end{figure}

\subsubsection{Limitations}

K-Means creates spherical clusters with hard boundaries, forcing every point into a
cluster regardless of whether it represents genuine attack behaviour or 
noise~\cite{sklearn-kmeans}. The PCA visualisation (Figure~\ref{fig:kmeans-pca}) 
highlights fragmentation: points assigned to the same cluster occupy widely separated 
regions of the feature space. These limitations prompted a transition to HDBSCAN.

\subsection{Algorithm 2: HDBSCAN Clustering}
\label{sec:hdbscan}

\subsubsection{Overview}

HDBSCAN~\cite{sklearn-hdbscan} (Hierarchical Density-Based Spatial Clustering of
Applications with Noise) addresses K-Means' limitations: it discovers clusters based
on density variations, automatically determines the number of clusters, identifies
noise points (labelled $-1$), and handles clusters of varying shapes and 
densities~\cite{hdbscan-params}.

The algorithm: (1) calculates core distances based on the $k$-th nearest neighbour;
(2) constructs a mutual-reachability distance graph; (3) builds a minimum spanning
tree; (4) converts it into a hierarchy of connected components (dendrogram); (5)
extracts stable clusters using the Excess of Mass method~\cite{sklearn-hdbscan}.

\begin{figure}[H]
  \centering
  \includegraphics[width=0.7\textwidth]{image93.png}
  \caption{Dendrogram Produced During HDBSCAN Hierarchical Clustering}
  \label{fig:dendrogram}
\end{figure}

\subsubsection{Advantages for Security Applications}

\begin{itemize}
  \item Explicitly identifies noise (label $-1$), distinguishing genuine attack
        patterns from random benign traffic~\cite{hdbscan-params}.
  \item Handles clusters of varying densities --- essential when different attack types
        exhibit vastly different activity levels~\cite{sklearn-hdbscan}.
  \item Multi-scale pattern detection via the cluster hierarchy captures reconnaissance
        followed by escalated exploitation as related 
        patterns~\cite{mitre-recon}.
  \item Parameters (\texttt{min\_cluster\_size}, \texttt{min\_samples}) have direct
        security interpretations: minimum attack footprint worth investigating and
        sensitivity to isolated suspicious events~\cite{hdbscan-params}.
\end{itemize}

\section{Model Training Results}

\subsection{Overview}

HDBSCAN produced substantially improved, analytically useful clusters compared to
K-Means~\cite{sklearn-hdbscan}. The UMAP projections show compact, well-separated 
clusters with clear behavioural patterns.

\begin{figure}[H]
  \centering
  \includegraphics[width=0.75\textwidth]{image94.png}
  \caption{3-D UMAP Projection of Trained HDBSCAN Model}
  \label{fig:umap-3d}
\end{figure}

\subsection{Cluster Labels and Analysis}
\label{sec:cluster-analysis}

\begin{table}[H]
  \centering
  \caption{HDBSCAN Cluster Labels with Confidence Scores (Retrained Model)}
  \label{tab:clusters}
  \small
  \begin{tabularx}{\textwidth}{cXcX}
    \toprule
    \textbf{Cluster} & \textbf{Label} & \textbf{Conf.} 
      & \textbf{Secondary Labels} \\
    \midrule
    $-1$ & PHP webshell deployment / backdoor probe      
      & 0.342 & PHPUnit RCE probing; Generic login discovery \\
      0  & Path traversal / directory escape probe       
      & 0.723 & --- \\
      1  & Git / source code exposure scan               
      & 0.774 & GeoServer / infrastructure admin probe; Secrets discovery \\
      2  & Generic login discovery / credential targeting 
      & 0.929 & Cisco / VPN / firewall probe; IoT exploit scan \\
      3  & Favicon / robots.txt fingerprinting           
      & 0.784 & PHP webshell deployment / backdoor probe \\
      4  & Favicon / robots.txt fingerprinting           
      & 0.820 & --- \\
      5  & Favicon / robots.txt fingerprinting           
      & 0.873 & --- \\
      6  & Uncategorized / miscellaneous scanning        
      & 0.000 & --- \\
      7  & PHP webshell deployment / backdoor probe      
      & 0.377 & Cloud metadata / config exposure; WordPress exploitation \\
      8  & PHP webshell deployment / backdoor probe      
      & 0.029 & Cloud metadata / config exposure \\
      9  & Exchange / Autodiscover / RDP probe           
      & 0.229 & Secrets discovery; Generic login discovery \\
     10  & PHP webshell deployment / backdoor probe      
      & 0.036 & --- \\
     11  & Exchange / Autodiscover / RDP probe           
      & 0.254 & Generic login discovery \\
     12  & PHP webshell deployment / backdoor probe      
      & 0.955 & WordPress exploitation / admin enumeration \\
     13  & Legitimate application traffic                
      & 0.598 & Secrets discovery; Cloud metadata / config exposure \\
     14  & IoT / embedded device exploit scan            
      & 0.220 & Exchange / Autodiscover probe; GeoServer probe \\
     15  & Uncategorized / miscellaneous scanning        
      & 0.000 & --- \\
     16  & IoT / embedded device exploit scan            
      & 0.202 & Cloud metadata / config exposure; Exchange probe \\
     17  & IoT / embedded device exploit scan            
      & 0.023 & Cisco / VPN / firewall probe \\
     18  & Secrets discovery (.env harvesting)           
      & 0.851 & PHP webshell probe; PHP info/config discovery \\
     19  & Secrets discovery (.env harvesting)           
      & 0.190 & Cisco / VPN probe; SonicWall SSLVPN probe \\
     20  & Secrets discovery (.env harvesting)           
      & 0.496 & PHP webshell probe; Cloud metadata / config exposure \\
     21  & WordPress exploitation / admin enumeration    
      & 1.000 & PHP webshell probe; ACME challenge path abuse \\
    \bottomrule
  \end{tabularx}
\end{table}

\subsubsection{Attack Distribution Summary}

Retraining on the expanded dataset produced 22 clusters (plus a noise group),
reflecting a richer and more granular threat landscape than the initial model:

\begin{enumerate}[label=\alph*.]
  \item \textbf{PHP Webshell / Backdoor Probing} (clusters $-1$, 7, 8, 10, 12):
        The dominant attack category, with attackers scanning for accessible PHP
        backdoors to establish persistent access~\cite{exploit-db}.
  \item \textbf{Secrets Discovery --- .env Harvesting} (clusters 18, 19, 20):
        Systematic attempts to retrieve environment files containing API keys,
        database credentials, and application 
        secrets~\cite{owasp-api8,github-sensitive-data}.
  \item \textbf{IoT / Embedded Device Exploitation} (clusters 14, 16, 17):
        Exploitation attempts targeting router, camera, and embedded device admin
        interfaces (HNAP, GPON, Hikvision, goform endpoints)~\cite{mitre-recon}.
  \item \textbf{Exchange / Autodiscover / RDP Probing} (clusters 9, 11):
        Scanning for exposed Microsoft Exchange, OWA, and RDP services a common
        ransomware initial-access vector~\cite{mitre-recon}.
  \item \textbf{Generic Login / Credential Targeting} (cluster 2):
        Broad scanning for login panels across multiple 
        frameworks~\cite{fortinet-brute-force}.
  \item \textbf{Favicon / robots.txt Fingerprinting} (clusters 3, 4, 5):
        Passive reconnaissance to identify technology stack before active 
        exploitation~\cite{mitre-recon}.
  \item \textbf{Git / Source Code Exposure} (cluster 1):
        Scanning for exposed \texttt{.git} directories that could leak full source
        code and embedded credentials~\cite{github-sensitive-data}.
  \item \textbf{Path Traversal / Directory Escape} (cluster 0):
        Attempts to escape the web root and access sensitive system files such as
        \texttt{/etc/passwd}~\cite{owasp-waf}.
  \item \textbf{WordPress Exploitation} (cluster 21):
        Enumeration of WordPress admin paths and exploitation of WP core 
        files~\cite{wp-king-addons}.
  \item \textbf{Legitimate Traffic} (cluster 13):
        Normal application traffic; suppressed from alerting unless anomalous
        signals appear~\cite{cf-verified-bots}.
  \item \textbf{Uncategorized / Miscellaneous} (clusters 6, 15):
        Scanning activity that did not match any known attack pattern.
\end{enumerate}

\subsubsection{Severity Classification}

\begin{table}[H]
  \centering
  \caption{Cluster Severity Ratings (Retrained Model)}
  \label{tab:severity}
  \small
  \begin{tabularx}{\textwidth}{cXcX}
    \toprule
    \textbf{Cluster} & \textbf{Label} & \textbf{Conf.} & \textbf{Severity} \\
    \midrule
    $-1$ & PHP webshell / backdoor probe      & 0.342 & \textbf{Critical} \\
      0  & Path traversal / directory escape  & 0.723 & \textbf{Critical} \\
      7  & PHP webshell / backdoor probe      & 0.377 & \textbf{Critical} \\
      8  & PHP webshell / backdoor probe      & 0.029 & \textbf{Critical} \\
     10  & PHP webshell / backdoor probe      & 0.036 & \textbf{Critical} \\
     12  & PHP webshell / backdoor probe      & 0.955 & \textbf{Critical} \\
      1  & Git / source code exposure scan    & 0.774 & High \\
      9  & Exchange / Autodiscover / RDP probe & 0.229 & High \\
     11  & Exchange / Autodiscover / RDP probe & 0.254 & High \\
     14  & IoT / embedded device exploit scan & 0.220 & High \\
     16  & IoT / embedded device exploit scan & 0.202 & High \\
     17  & IoT / embedded device exploit scan & 0.023 & High \\
     18  & Secrets discovery (.env harvesting) & 0.851 & High \\
     19  & Secrets discovery (.env harvesting) & 0.190 & High \\
     20  & Secrets discovery (.env harvesting) & 0.496 & High \\
     21  & WordPress exploitation              & 1.000 & High \\
      2  & Generic login / credential targeting & 0.929 & Medium \\
      3  & Favicon / robots.txt fingerprinting & 0.784 & Low \\
      4  & Favicon / robots.txt fingerprinting & 0.820 & Low \\
      5  & Favicon / robots.txt fingerprinting & 0.873 & Low \\
      6  & Uncategorized / miscellaneous       & 0.000 & Low \\
     15  & Uncategorized / miscellaneous       & 0.000 & Low \\
     13  & Legitimate application traffic      & 0.598 & Info \\
    \bottomrule
  \end{tabularx}
\end{table}

\section{Model Testing and Drift Analysis}
\label{sec:drift}

After fitting and fine-tuning, the retrained model was evaluated against a withheld
test dataset. Performance was measured by the \textit{cluster drift score}: the
Euclidean distance between training and test cluster centroids in UMAP 
space~\cite{concept-drift}. Higher scores indicate greater behavioural drift or novel 
attack patterns in the test window; lower scores confirm that a cluster represents a 
stable, repeatable tactic.

\begin{table}[H]
  \centering
  \caption{Cluster Drift Scores After Testing (Retrained Model, sorted by drift)}
  \label{tab:drift}
  \begin{tabular}{lSSSS}
    \toprule
    \textbf{Cluster} & \textbf{UMAP x} & \textbf{UMAP y} & \textbf{UMAP z} 
      & \textbf{Drift Score} \\
    \midrule
     4   &  44.957128 & 11.870636 &  6.747134 & 14.526147 \\
     2   &  24.913916 &  5.581988 &  0.859545 &  7.486253 \\
    11   &  11.897444 &  0.345456 &  2.958475 &  3.531601 \\
    13   &   0.031693 &  0.679325 &  1.656599 &  1.790756 \\
    20   &   0.106773 &  1.697355 &  0.416053 &  1.750861 \\
    $-1$ &   0.470403 &  1.612247 &  0.024043 &  1.679642 \\
    20   &   0.300752 &  1.401953 &  0.426247 &  1.495865 \\
     5   &   0.792748 &  0.599638 &  0.949964 &  1.374936 \\
    18   &   0.612942 &  0.478292 &  0.624942 &  0.997503 \\
    19   &   0.142550 &  0.102555 &  0.855991 &  0.873819 \\
    12   &   0.171323 &  0.406830 &  0.122215 &  0.458038 \\
    14   &   0.400525 &  0.018703 &  0.192242 &  0.444665 \\
     3   &   0.040725 &  0.376230 &  0.215489 &  0.435480 \\
    16   &   0.297522 &  0.170392 &  0.132258 &  0.367484 \\
     6   &   0.248292 &  0.194670 &  0.021723 &  0.316255 \\
    11   &   0.048236 &  0.224376 &  0.206218 &  0.308540 \\
     7   &   0.219868 &  0.027002 &  0.178785 &  0.284667 \\
    17   &   0.148715 &  0.235938 &  0.030113 &  0.280517 \\
    15   &   0.011355 &  0.162050 &  0.159172 &  0.227431 \\
    10   &   0.063951 &  0.085261 &  0.078014 &  0.132082 \\
    \bottomrule
  \end{tabular}
\end{table}

Clusters 4 and 2 exhibit the highest drift scores (14.53 and 7.49 respectively),
suggesting that favicon fingerprinting and generic login-targeting traffic shifted
the most between training and test windows likely reflecting the opportunistic,
tool-driven nature of these scans where attackers rotate tooling 
frequently~\cite{mitre-recon}. Cluster $-1$ (noise/webshell probing) shows a moderate 
drift of 1.68, consistent with evolving backdoor payload lists~\cite{exploit-db}. The 
majority of clusters, particularly the stable attack types such as .env harvesting 
(cluster 18, drift 1.00), IoT probing (cluster 14, drift 0.44), and WordPress 
exploitation (cluster 21, not represented in test window), show low drift scores 
confirming that these represent persistent, repeatable tactics rather than one-off 
anomalies~\cite{concept-drift}. Monitoring drift scores over time provides a practical 
trigger for model retraining when scores begin to climb consistently across multiple 
clusters~\cite{concept-drift}.

\begin{figure}[H]
  \centering
  \includegraphics[width=0.7\textwidth]{image96.png}
  \caption{UMAP Visualisation of Cluster Drift After Testing 
  (Blue = Training, Red = Test)}
  \label{fig:drift-umap}
\end{figure}

% ============================================================
%  CHAPTER 6 — IMPLEMENTATION
% ============================================================
\chapter{Implementation}
\label{chap:implementation}

\section{Overview}

The Jupyter Notebook pipeline was converted into two production Python scripts:
\begin{itemize}
  \item \href{https://github.com/G-cae78/ProtectTheHoney/blob/main/Machine\%20Learning/GCP-VM-Pipeline\%20/run_batch.py}{\texttt{run\_batch.py}}
  \item \href{https://github.com/G-cae78/ProtectTheHoney/blob/main/Machine\%20Learning/GCP-VM-Pipeline\%20/pipeline.py}{\texttt{pipeline.py}}
\end{itemize}

Dependencies were packaged and deployed to the GCP virtual machine. A cron job
triggers the clustering pipeline~\cite{sklearn-hdbscan} every 5 minutes on the most 
recent web traffic logs.

\begin{figure}[H]
  \centering
  \includegraphics[width=\textwidth]{image97.png}
  \caption{Prediction Summary and Rollup Output on GCP Instance}
  \label{fig:pipeline-output}
\end{figure}

\section{Pipeline Deployment}

\subsection{Output Files}

Each pipeline execution produces two outputs:

\paragraph{rollup.json}
Contains cluster assignments for all traffic observed in the previous 5-minute
period, including top targeted paths, cluster strength, and noise 
rate~\cite{sklearn-hdbscan}.

\begin{figure}[H]
  \centering
  \includegraphics[width=\textwidth]{image98.png}
  \caption{rollup.json File Structure}
  \label{fig:rollup-json}
\end{figure}

\paragraph{pred.csv}
Summarises the IP addresses active during that period, their recommended security
action level (Low, Medium, High, or Critical), and aggregate request counts per IP
and cluster~\cite{nist-800-53}. Table~\ref{tab:pred-csv} shows an example record.

\begin{table}[H]
  \centering
  \caption{Example Record from pred.csv}
  \label{tab:pred-csv}
  \begin{tabular}{ll}
    \toprule
    \textbf{Field} & \textbf{Value} \\
    \midrule
    Client IP          & 193.142.147.209 \\
    Window Start (UTC) & 2026-01-30 15:35:00 \\
    Request Count      & 1 \\
    Unique Paths       & 1 \\
    Unique Methods     & 1 \\
    Status 200         & 0 \\
    Status 401         & 0 \\
    Status 403         & 0 \\
    Status 404         & 0 \\
    Status 429         & 0 \\
    4xx Ratio          & 1.000 \\
    Path               & / \\
    Cluster            & 4 \\
    Cluster Strength   & 0.999401 \\
    User Agents        & Mozilla/5.0 (Windows NT 10.0; Win64; x64) \ldots \\
    \bottomrule
  \end{tabular}
\end{table}

\section{Job Scheduling and Monitoring}

Initially, pipeline jobs were triggered using standard cron jobs executed directly
from the instance terminal. While functional, this approach offered no visibility
into run history, no alerting on failure, and no resource tracking. To improve
monitoring and reliability, scheduling was migrated to
Cronicle~\cite{cronicle}, a web-based job scheduler and monitoring system.

Cronicle provides a visual interface showing the status of each job run with clear
indications of successful completions and failures. The platform sends notifications
when jobs fail, enabling quicker troubleshooting and resolution if something goes
wrong. It also tracks performance metrics per execution, CPU usage, memory
consumption, and wall-clock execution time which have been valuable in confirming
the pipeline does not approach the VM's resource limits~\cite{cronicle}. Based on 
the metrics collected across hundreds of runs, it is far from doing so: each run 
averages 2 seconds and consumes only 0.5\% of CPU.

\begin{figure}[H]
  \centering
  \includegraphics[width=0.85\textwidth]{image99.png}
  \caption{Cronicle Job Completion Dashboard Showing Resource Usage Per Pipeline Run}
  \label{fig:cronicle-dash}
\end{figure}

\begin{figure}[H]
  \centering
  \includegraphics[width=0.57\textwidth]{image100.png}
  \caption{GCP Instance Utilisation (GCP Observability Tab)}
  \label{fig:gcp-util}
\end{figure}

Figure~\ref{fig:cronicle-dash} shows the Cronicle interface for a successful job
completion. The job executed in 3 seconds, consuming 0\% of CPU and negligible
memory. The success banner confirms the clustering pipeline completed without errors,
and the blue CPU bar indicates that all processing on the instance during that window
was the pipeline script itself, with minimal system overhead. This demonstrates the
pipeline's efficiency: it is lightweight enough to run every 5 minutes without
impacting any other processes on the VM.

\begin{figure}[H]
  \centering
  \includegraphics[width=\textwidth]{image101.png}
  \caption{Pipeline Job History in Cronicle}
  \label{fig:cronicle-history}
\end{figure}

The completed-jobs tab highlights the reliability of the pipeline by listing all
historical runs with timestamps, each separated by the expected 5-minute interval.
Uniform success statuses across all runs confirm the deployment is stable and does
not encounter errors during normal operation. This view also functions as an audit
trail: analysts can verify the clustering model ran continuously during any given
time period, providing far better forensic visibility than a raw cron log in a
terminal~\cite{aws-cloudwatch}.

In Cronicle's forward-looking schedule view, each row represents a forthcoming 
execution, displaying the exact command being run (including the Python script path 
\texttt{/opt/honey/pipeline/run\_batch\_nginx.py}), the plugin type (shell script), 
and a countdown to the next trigger~\cite{cronicle}. The \textit{Category} column 
identifies these as general jobs while the \textit{Plugin} column confirms they are 
shell-script executions of the Python clustering pipeline. Cronicle manages the 
recurring schedule without drift or gaps, and the \textit{Edit Event} links allow 
individual runs to be adjusted or paused for maintenance without disrupting the rest 
of the queue.

Overall, Cronicle~\cite{cronicle} gives the project a production-like feel with 
comprehensive functionality for maintenance, testing, error notifications, and 
telemetry. The platform transforms what would have been an invisible cron job into a 
fully observable, manageable system with real-time status monitoring, historical audit
trails, and proactive failure alerting. The performance metrics collected over
hundreds of successful runs provide confidence that the clustering pipeline is both
reliable and efficient. The combination of automated execution with comprehensive
monitoring demonstrates that the clustering-based threat detection system can
function as a robust, enterprise-ready solution.

\section{Agentic LLM \& Interpretability}
\label{sec:agentic-llm}

\subsection{Overview}

While the clustering model assigns traffic windows to behavioural clusters and scores
confidence, a label like ``IoT / embedded device exploit scan'' with a confidence of
0.023 provides little actionable guidance~\cite{d3fend}. The LLM layer does not 
replace the clustering output; it interprets and enriches it into something actionable.

\begin{figure}[H]
  \centering
  \includegraphics[width=\textwidth]{image103.png}
  \caption{Agentic LLM Pipeline Architecture}
  \label{fig:llm-arch}
\end{figure}

\subsection{Context Supplied to the Model}

On each pipeline run the model receives:

\begin{itemize}
  \item The current rollup and prediction rows.
  \item Cluster severity labels~\cite{nist-800-53}.
  \item Active Cloudflare WAF rules as context~\cite{cf-api-lockdown}.
\end{itemize}

From this it produces a human-readable threat summary, specific WAF actions tied to
each detected cluster, and a triage priority list ordered by severity. For every
Medium, High, or Critical cluster it raises a structured Jira ticket~\cite{jira-api} 
containing the observed activity, matched patterns, recommended action, and supporting 
references. The model is explicitly instructed not to fabricate CVE URLs to prevent
hallucinations.

\subsection{RAG and Knowledgebase Integration}

The MITRE D3FEND~\cite{d3fend} knowledgebase (thousands of entries) cannot be
pasted wholesale into the prompt. Instead, each entry is pre-embedded offline using
\texttt{voyageai/voyage-3-lite} and stored as normalised vectors. At query time, the 
model calls a \texttt{query\_knowledgebase} tool, and cosine-similarity retrieval 
injects the three most relevant entries into the prompt, keeping the context window 
lean while grounding recommendations in established 
frameworks~\cite{d3fend,mitre-recon}.

\subsection{Agentic Tool-Use Loop}

The model operates as a tool-use loop, making sequential calls within a single run:

\begin{enumerate}
  \item Query the MITRE D3FEND knowledgebase~\cite{d3fend} once per detected cluster 
        to retrieve severity labels, matched patterns, and threat intelligence to 
        guide remediation.
  \item Call \texttt{create\_consolidated\_firewall\_rules} to deploy or update
        Cloudflare WAF rules via the Rulesets API~\cite{cf-api-lockdown}, 
        consolidating all detected clusters into a minimum number of OR-chained 
        expressions to respect the free-tier five-rule budget.
  \item Call \texttt{create\_jira\_ticket}~\cite{jira-api} for every Medium, High, 
        or Critical cluster, embedding the rule ID, action taken, and knowledgebase 
        references directly into the ticket description so engineers have full context 
        without consulting external logs.
\end{enumerate}

Cluster 13 (known legitimate application traffic, labelled ``Legitimate application
traffic'' with confidence 0.598) is suppressed from ticket creation unless anomalous
signals appear, such as secrets-harvesting paths~\cite{owasp-api8} or an elevated 
4xx ratio.

\subsection{Deployment and Scheduling}

The pipeline was implemented as a FastAPI application deployed on the same GCP
Compute Engine instance as the clustering model. A \texttt{pipeline\_sender.py}
script monitors the \texttt{/opt/honey/results} directory and batches every two
rollups (ten minutes of traffic) into a single HTTP POST to the
\texttt{/pipeline/rollup/batch} endpoint. Batching reduces the number of Groq API
calls and avoids per-request overhead; three to five rollups are processed in a
single request. Both services are managed by \texttt{systemd} and restart
automatically on failure or reboot.

\subsection{Embedding Cache and Token Optimisation}

Pre-computing embeddings offline eliminated the dominant token cost. The MITRE
D3FEND knowledgebase~\cite{d3fend}, supplemented with cluster severity metadata and 
Cloudflare documentation pages~\cite{cf-how-works}, was embedded using 
\texttt{voyage-3-lite} and persisted to disk as a \texttt{.pkl} file via Python's 
\texttt{pickle} library. The resulting index contains 4,668 entries of dimension 512. 
A MD5 hash of the full text corpus is stored alongside the vectors; if the corpus 
changes the cache is automatically invalidated and re-embedded. At runtime the file 
is loaded once on startup and held in memory, reducing per-cycle token usage by 70.4\% 
compared to injecting the full knowledgebase as a system message.

\begin{table}[H]
\centering
\caption{Token usage before and after embedding optimisation}
\begin{tabular}{lrr}
\hline
\textbf{Component} & \textbf{Before} & \textbf{After} \\
\hline
HDBSCAN cluster context   & 2,599 & 23    \\
Cluster severity levels   & 2,202 & 0     \\
Current predictions       & 160   & 0     \\
Current rollup            & 164   & 164   \\
Prompt                    & 1,889 & 1,889 \\
\hline
\textbf{Total per cycle}  & \textbf{7,014} & \textbf{2,076} \\
\hline
\multicolumn{3}{l}{Reduction: 70.4\% \quad Savings per day (288 cycles): 
  1,422,144 tokens} \\
\hline
\end{tabular}
\label{tab:token-optimisation}
\end{table}

\subsection{Model Availability and Rate Limit Mitigation}

Groq's free tier imposes per-model daily token limits (TPD). To guarantee full
24-hour coverage across all 288 five-minute cycles, a fallback chain of seven models
is configured~\cite{qwen3,hf-inference}. The pipeline attempts each model in priority 
order, falling back automatically on rate-limit (HTTP 429), server error (5xx), or 
tool-calling failure responses.

\begin{table}[H]
\centering
\caption{Groq model fallback chain and daily cycle capacity}
\begin{tabular}{lrr}
\hline
\textbf{Model} & \textbf{TPD} & \textbf{Max cycles} \\
\hline
qwen/qwen3-32b (primary)                      & 500,000   & 139 \\
meta-llama/llama-4-scout-17b-16e-instruct     & 500,000   & 139 \\
llama-3.1-8b-instant                          & 500,000   & 139 \\
moonshotai/kimi-k2-instruct                   & 300,000   & 83  \\
openai/gpt-oss-20b                            & 200,000   & 55  \\
openai/gpt-oss-120b                           & 200,000   & 55  \\
llama-3.3-70b-versatile                       & 100,000   & 27  \\
\hline
\textbf{Total combined} & \textbf{2,300,000} & \textbf{643} \\
\hline
\multicolumn{3}{l}{Cycles needed per day: 288 \quad Coverage: \checkmark Full day} \\
\hline
\end{tabular}
\label{tab:model-fallback}
\end{table}

The combined TPD of 2,300,000 tokens across all models yields 643 available cycles
against the 288 required, providing a 2.2$\times$ safety margin~\cite{qwen3}. 
\texttt{qwen/qwen3-32b} serves as the primary model due to consistent tool-calling 
reliability observed during testing~\cite{qwen3}. All selected models support the 
Groq tool-calling API and have been verified to produce valid structured outputs for 
the three pipeline tools~\cite{hf-inference}.
```

% ============================================================
%  CHAPTER 7 — RESULTS  (placeholder)
% ============================================================
\chapter{Results and Discussion}
\label{chap:results}

This chapter synthesises findings across the CDN comparative 
evaluation (Chapter~\ref{chap:experimentation}), the HDBSCAN 
traffic clustering model (Section~\ref{sec:hdbscan}), and 
the agentic LLM interpretability pipeline 
(Section~\ref{sec:agentic-llm}), drawing conclusions against 
each research question stated in 
Section~\ref{sec:research-questions}.

\section{CDN Comparative Findings}
\label{sec:cdn-findings}

The full comparative results are presented in 
Table~\ref{tab:comparison} (Section~\ref{sec:comparison}). The 
findings reveal a consistent pattern: both platforms handled 
network-layer attacks reliably, but diverged significantly at 
the application layer where most real-world threats operate.

At Layer 3 and 4, both CDNs absorbed SYN floods, UDP floods, 
and ICMP floods entirely at the edge with zero measurable impact 
on the origin server. Neither platform is meaningfully 
differentiated at this layer, and for organisations primarily 
concerned with volumetric attacks, the free tier of either 
provider is sufficient.

The gap opens at Layer 7. Cloudflare's free tier allowed 
31--60\% of HTTP/2 flood traffic to reach the origin across 
repeated test runs, with inconsistent mitigation between 
identical runs suggesting the detection threshold is 
non-deterministic~\cite{http2-ddos}. BunnyCDN escalated 
progressively to a full IP-level ban by the fourth run, 
reaching zero origin impact~\cite{bunny-shield-docs}. For slow 
HTTP attacks, Cloudflare mitigated silently with no dashboard 
visibility, while BunnyCDN logged 293 bot-detection events, 
giving an administrator at least the awareness that an attack 
had occurred.

The most operationally significant finding is the SQL injection 
disparity. Across four test runs using SQLMap at maximum 
aggression, tamper scripts, and a dedicated WAF-bypass 
framework, Cloudflare's free tier generated zero WAF events and 
applied zero mitigations. A single \texttt{--random-agent} flag 
was sufficient to make 59,950 requests appear entirely benign. 
BunnyCDN detected and logged every run, with its libjection 
module normalising obfuscated payloads before matching, meaning 
tamper scripts provided no evasion benefit. SQL injection 
remains consistently ranked among the most critical web 
application vulnerabilities~\cite{owasp-waf}, and for any 
application not using parameterised queries, Cloudflare's free 
tier offers no meaningful protection out of the box.

These findings suggest the security gap between the two 
platforms is not primarily a function of price but of default 
rule depth. BunnyCDN Advanced Shield ships with 255 base WAF 
rules and a libjection normalisation layer active by 
default~\cite{bunny-waf-docs}. Cloudflare's free tier relies 
more heavily on heuristics and user-agent matching, which are 
trivially bypassed~\cite{cf-how-works}. Organisations choosing 
a CDN on security grounds should treat application-layer WAF 
coverage as the critical differentiator.

\section{Closing the Gap with Automated Rule Consolidation}
\label{sec:rule-consolidation}

The CDN comparison above raises an obvious question: if 
Cloudflare's free tier ships with weak default WAF coverage, 
can supplementary custom rules close that gap? The answer from 
this project is yes, but only if those rules are generated 
intelligently and maintained automatically.

Cloudflare's free tier permits five custom WAF rules per zone. 
On the surface this appears severely limiting, but the 
expression language imposes no limit on logical operators within 
a single rule, only a 4,096-character ceiling on the total 
expression string~\cite{cf-community-blocking}. This means a 
single rule can encode twenty or more distinct threat signatures 
using OR-chained path patterns, effectively functioning as a 
sub-rule container. The agentic pipeline exploits this 
directly: rather than consuming one budget slot per detected 
cluster, it groups clusters of equivalent severity and action 
tier into shared expressions, appending new patterns as they 
are detected without consuming additional slots.

The critical distinction is that the pipeline does not generate 
arbitrary rules. Each rule expression is derived from the 
specific path patterns observed in a named HDBSCAN 
cluster~\cite{sklearn-hdbscan}, grounded in MITRE D3FEND 
knowledgebase entries retrieved at query time~\cite{d3fend}. 
A rule blocking \texttt{/\%2eenv} variants is deployed because 
the clustering model identified a statistically coherent 
secrets-harvesting behaviour pattern, not because a generic 
blocklist was applied. This means the five-rule budget is spent 
on threats actually observed against this specific deployment, 
rather than on broad signatures that may generate false 
positives against legitimate traffic.

Over the three-day evaluation period, seven semantically 
distinct threat categories were covered within the five-rule 
budget, intercepting 999 events with only five confirmed false 
positives, all corrected by a single prompt adjustment. This 
demonstrates that the constraint of five rules is not a ceiling 
on protection coverage but on the number of independently 
manageable expressions, and that an LLM with access to cluster 
context can saturate that budget with high-precision, 
deployment-specific signatures that a human analyst would take 
significantly longer to write and maintain manually.

\section{Agentic Pipeline Evaluation}
\label{sec:pipeline-results}

\subsection{WAF Rule Effectiveness}
To evaluate the end-to-end pipeline, three days of Cloudflare 
firewall event logs were exported covering the period 19--22 
March 2026. The logs capture every request matched by a rule 
created autonomously by the LLM pipeline, with no manual 
intervention at any stage after deployment~\cite{cf-api-lockdown}.

Across the three-day window the pipeline-created and adjusted 
rules intercepted \textbf{999 unique events}, broken down as 
follows:

\begin{table}[H]
\centering
\caption{WAF rule actions over the evaluation period 
(19--22 March 2026)}
\begin{tabular}{lrr}
\hline
\textbf{Action} & \textbf{Events} & \textbf{Proportion} \\
\hline
Interactive challenge  & 913 & 91.4\% \\
Hard block             &  67 &  6.7\% \\
Managed challenge      &  19 &  1.9\% \\
\hline
\textbf{Total}         & \textbf{999} & \textbf{100\%} \\
\hline
\end{tabular}
\label{tab:waf-actions}
\end{table}

The distribution reflects the pipeline's tiered severity 
strategy: Critical-severity clusters map to \texttt{block}, 
High-severity clusters to \texttt{managed\_challenge}, and 
Medium-severity clusters to \texttt{challenge}. The dominance 
of challenge events indicates the pipeline correctly identified 
the majority of threats as automated scanning rather than 
targeted exploitation, preserving the possibility of legitimate 
access while impeding bot traffic~\cite{cf-bot-fight}.

\subsection{Rule Coverage and Consolidation}
As described in Section~\ref{sec:rule-consolidation}, 
Cloudflare's five-rule limit is not a ceiling on protection 
coverage when rules are OR-chained intelligently. Over the 
evaluation period the pipeline deployed the following rules 
autonomously:

\begin{table}[H]
\centering
\caption{Pipeline-created WAF rules and event counts}
\begin{tabular}{p{8.5cm}rr}
\hline
\textbf{Rule description} & \textbf{Events} & \textbf{Action}\\
\hline
Medium severity cluster 2 --- IoT admin login scan        
  & 482 & challenge \\
Challenge zgrab user-agent OWA endpoint probing           
  & 385 & challenge \\
CRITICAL: PHP webshell / WordPress admin enumeration      
  &  53 & block \\
Block low severity cluster 1 scanning                     
  &  40 & block \\
HIGH: IoT/embedded device exploit scan (cluster 17)       
  &  17 & managed\_challenge \\
CRITICAL: Cluster 17 --- unknown/noise root path activity 
  &  14 & block \\
Low: Unclassified traffic pattern (cluster 10)            
  &   5 & challenge \\
\hline
\end{tabular}
\label{tab:rule-events}
\end{table}

Seven semantically distinct threat categories are handled 
within the five-rule budget by grouping clusters of equivalent 
severity and action tier into shared OR-chained expressions. 
As new clusters appear in subsequent rollups the pipeline 
appends their patterns to the appropriate existing rule rather 
than consuming an additional budget slot, allowing protection 
coverage to grow continuously without manual intervention.

\subsection{Persistent Threat Actor: BUCKLOG}
The most significant finding across the evaluation period was 
a persistent automated scanning campaign originating from ASN 
\texttt{BUCKLOG} (France), accounting for \textbf{874 of the 
999 events (87.5\%)}. The actor operated from only three 
distinct IP addresses and almost exclusively used 
\texttt{curl/8.7.1} as its user agent, confirming fully 
automated tooling. Two coordinated bursts were observed at 
06:00 UTC on consecutive days (444 and 391 events 
respectively), consistent with a scheduled cron-driven scan.

The targeted paths reveal a sophisticated evasion strategy, 
consistent with documented secrets-harvesting and environment 
file exposure techniques~\cite{owasp-api8,github-sensitive-data}:

\begin{itemize}
  \item \texttt{/\%2eenv} --- URL-encoded \texttt{.env} to 
        bypass naive string matching on the literal dot 
        character.
  \item \texttt{/\$(pwd)/.env.staging} --- shell metacharacter 
        injection attempting to resolve environment files 
        relative to the working directory.
  \item \texttt{/etc/grafana/grafana.ini} --- targeting Grafana 
        credential storage.
  \item \texttt{/api/v1/credentials}, \texttt{/api/credentials} 
        --- systematic enumeration of credential API endpoints.
  \item \texttt{/var/www/.env}, \texttt{/stripe/.env} --- 
        directory-prefix variants of secrets harvesting.
\end{itemize}

When the pipeline first detected the BUCKLOG cluster it 
followed its tiered severity strategy and issued a challenge 
response, the appropriate initial action for a pattern that 
had not yet demonstrated persistence. Cloudflare's free tier, 
left to its own default rules, had no equivalent response at 
all: the same actor and the same paths generated zero WAF 
events natively, as established in the SQL injection testing 
in Section~\ref{sec:cf-sqli}. The pipeline at least 
intercepted the traffic from the first detection cycle.

The challenge proved insufficient. A fully automated tool 
with no human in the loop absorbs challenge overhead at 
negligible cost~\cite{cf-captcha-end}, and BUCKLOG returned 
the following day with an identical burst at the same time. 
On detecting the recurrence across successive rollup cycles, 
the pipeline escalated its response, promoting the action tier 
from challenge to block for the BUCKLOG-associated patterns. 
This escalation happened autonomously, without any manual 
intervention, and is the key operational result of the 
evaluation: not just that the pipeline detected the threat, 
but that it recognised the challenge was not working and 
changed its strategy accordingly.

All 874 events were ultimately challenged or blocked despite 
the URL-encoding and path-manipulation evasion techniques. 
This is further notable because the HDBSCAN model was trained 
on unencoded paths; the LLM's use of 
\texttt{http.request.uri.path contains} expressions in 
Cloudflare's normalised field, which decodes percent-encoding 
before matching, provided implicit coverage of encoded 
variants without any explicit training on them. Bare 
Cloudflare, without the pipeline, would have allowed every 
one of these requests through silently.

\subsection{Temporal Attack Pattern}
Figure~\ref{fig:waf-timeline} shows the hourly event 
distribution across the evaluation window. The two 06:00 UTC 
spikes are separated by exactly 24 hours, strongly suggesting 
a scheduled scanning job~\cite{mitre-recon}. The inter-burst 
traffic represents background noise from opportunistic 
scanners on other ASNs including Microsoft, Google Cloud, 
and Tencent.

\begin{figure}[H]
  \centering
  \includegraphics[width=\textwidth]{waf_timeline.png}
  \caption{Hourly WAF rule event distribution, 
  19--22 March 2026}
  \label{fig:waf-timeline}
\end{figure}

\subsection{Geographic and Network Distribution}

\begin{table}[H]
\centering
\caption{Top source countries by event count}
\begin{tabular}{lrr}
\hline
\textbf{Country} & \textbf{Events} & \textbf{Proportion} \\
\hline
France (FR)        & 877 & 87.8\% \\
United States (US) &  76 &  7.6\% \\
Lithuania (LT)     &   9 &  0.9\% \\
Germany (DE)       &   7 &  0.7\% \\
China (CN)         &   5 &  0.5\% \\
Ireland (IE)       &   5 &  0.5\% \\
Other              &  20 &  2.0\% \\
\hline
\end{tabular}
\label{tab:geo-distribution}
\end{table}

The concentration in France is attributable entirely to the 
BUCKLOG campaign. United States traffic originates 
predominantly from cloud infrastructure and represents a mix 
of legitimate crawlers~\cite{cf-verified-bots} and 
opportunistic scanners caught by broad challenge rules. The 
Irish traffic is consistent with the honeypot's European 
hosting location attracting regional scanners.

\subsection{Pipeline Autonomy Assessment}
Across the evaluation period the pipeline operated without any 
manual intervention. Starting from raw HDBSCAN cluster 
rollups, it autonomously:

\begin{enumerate}
  \item Queried the MITRE D3FEND knowledgebase~\cite{d3fend} 
        to ground recommendations in established threat 
        frameworks.
  \item Deployed and iteratively updated Cloudflare WAF rules 
        via the Rulesets API~\cite{cf-api-lockdown}, 
        OR-chaining new cluster patterns into existing rules 
        as they were detected.
  \item Raised structured Jira tickets~\cite{jira-api} for 
        every Medium, High, and Critical cluster, each 
        containing the rule ID, patterns deployed, and 
        knowledgebase references.
  \item Adapted rule expressions to cover URL-encoded and 
        path-manipulated variants of known attack patterns 
        through Cloudflare's normalised field matching.
\end{enumerate}

The results demonstrate that an unsupervised clustering 
model~\cite{sklearn-hdbscan} combined with an agentic LLM 
layer can replace the manual analyst workflow of identifying 
threats, writing WAF expressions, and filing incident tickets, 
all within the constraints of a free-tier infrastructure. The 
system processed 288 five-minute rollup cycles per day across 
the evaluation period, with the model fallback chain ensuring 
uninterrupted coverage despite individual model rate limits 
on Groq's free tier.

\subsection{Limitations}
Several limitations were observed during evaluation. The model 
occasionally mapped Critical-severity clusters to 
\texttt{managed\_challenge} rather than \texttt{block}, 
reflecting imprecise severity-to-action mapping in the prompt 
rather than a failure of cluster detection. The pipeline also 
depends on Groq's free-tier availability; while the 
seven-model fallback chain provides a 2.2$\times$ daily 
capacity margin, a sustained outage across all models would 
interrupt rule deployment until the next cycle. Finally, 
Cloudflare's free tier does not expose Security Events beyond 
a 24-hour window~\cite{cf-ddos}, requiring an export script 
to run at least twice daily to maintain a complete historical 
record.

% ============================================================
%  CHAPTER 8 — FUTURE WORK
% ============================================================
\chapter{Future Work}
\label{sec:future-work}

\section{Multi-CDN Comparative Testing}

The agentic pipeline is currently deployed behind Cloudflare's 
free tier, which imposes hard constraints that directly limit 
what the pipeline can do: five custom WAF rules per zone and 
minimal telemetry for network-layer events. These constraints 
forced the pipeline to OR-chain multiple threat signatures into 
single rules rather than deploying them independently, reducing 
granularity and making individual rules harder to audit.

Other free-tier providers such as AWS CloudFront with WAF and 
Google Cloud Armor have different rule budgets, logging 
capabilities, and expression languages. Pairing the same 
agentic pipeline with each provider would reveal which free-tier 
offering gives the pipeline the most room to operate, and whether 
a more permissive rule budget produces meaningfully better 
detection coverage. It is possible that a combination approach, 
running the pipeline behind multiple CDNs simultaneously, could 
close the gap to paid-tier protection without any additional 
spend, since each provider's strengths would compensate for 
another's limitations.

\section{Automated Drift Detection and Retraining Alerts}

Retraining the HDBSCAN model is currently triggered manually. 
A scheduled daily job could compute per-cluster drift scores 
automatically and raise an alert when any cluster's centroid 
displacement exceeds a configurable threshold, signalling that 
a known attack pattern has evolved enough to warrant retraining. 
Doing this at the per-cluster level means a single shifting 
category can be flagged without affecting well-settled clusters 
elsewhere, keeping the model current without unnecessary 
disruption.

\section{Reinforcement Learning for Adaptive Threat Detection}

Rather than retraining a static model periodically, a 
reinforcement learning agent would update its policy continuously 
from live traffic. It would receive positive reward when flagged 
clusters are confirmed as genuine attacks and negative reward 
when rules produce false positives, gradually learning when to 
issue a hard block versus a softer challenge. Combined with 
automated drift detection, this would produce a fully 
self-sustaining pipeline that improves from experience without 
ever needing to be taken offline for retraining.

\clearpage
\chapter{Sustainability Statement}
\addcontentsline{toc}{chapter}{Sustainability Statement}

This project was designed with sustainability considerations across 
environmental, social, and economic dimensions, and maps to three 
UN Sustainable Development Goals: SDG 9 (Industry, Innovation and 
Infrastructure), SDG 16 (Peace, Justice and Strong Institutions), 
and SDG 17 (Partnerships for the Goals).

\textbf{Environmental.} The honeypot ran on a single GCP 
\texttt{e2-micro} instance consuming roughly 0.2--0.5W at idle, 
totalling approximately 1--2 kWh across the six-month deployment. 
CDN edge caching reduced origin server load during high-traffic 
periods, and all LLM inference was offloaded to Groq's LPU-based 
API, removing the need for a continuously running GPU instance.

\textbf{Economic.} The full stack, covering clustering, visualisation, 
RAG, and agentic response, was built entirely on free and open-source 
tooling within cloud free-tier limits. This shows that a meaningful 
automated threat detection pipeline is achievable without 
enterprise-grade spending, directly lowering the barrier for 
smaller organisations to adopt similar approaches (SDG 9).

\textbf{Social.} The project contributes a replicable methodology 
for small teams to classify live attack traffic without labelled 
training data and automate incident response with freely available 
tools. Grounding recommendations in the MITRE D3FEND knowledgebase 
promotes community-maintained, vendor-neutral security standards, 
which is particularly valuable for organisations that cannot afford 
proprietary security frameworks (SDG 16, SDG 17).

\textbf{Limitations.} A continuously running VM incurs an ongoing 
carbon cost regardless of traffic volume, and the Cloudflare and 
Jira integrations introduce dependency on third-party commercial 
platforms. Future iterations could explore spot-instance or 
serverless deployments to reduce idle energy consumption.

\phantomsection
\addcontentsline{toc}{chapter}{References}
\renewcommand{\bibname}{References}
\begin{thebibliography}{99}

% ... your \bibitem entries ...


\bibitem{serverless-honeypot}
H. Varna, ``Building a Free Serverless Honeypot on AWS (Lambda + API Gateway),''
\textit{Medium}. [Online]. Available:
\url{https://hemavarnas.medium.com/build-a-free-serverless-honeypot-on-aws-lambda-api-gateway-e9b896d7fe54}

\bibitem{cf-reverse-proxy}
Cloudflare, ``What Is a Reverse Proxy?,'' \textit{Cloudflare Learning Center}.
[Online]. Available: \url{https://www.cloudflare.com/en-gb/learning/cdn/glossary/reverse-proxy/}

\bibitem{cf-how-works}
Cloudflare, ``How Cloudflare Works,'' \textit{Cloudflare Developer Docs}.
[Online]. Available: \url{https://developers.cloudflare.com/fundamentals/concepts/how-cloudflare-works/}

\bibitem{duckdns}
DuckDNS, ``Free Dynamic DNS Hosted on AWS,'' \textit{DuckDNS}. [Online]. Available:
\url{https://www.duckdns.org}

\bibitem{lovable}
Lovable AI, ``Build Something With Lovable,'' \textit{Lovable}. [Online]. Available:
\url{https://lovable.dev}

\bibitem{awesome-honeypots}
paralax, ``Awesome Honeypots,'' \textit{GitHub}. [Online]. Available:
\url{https://github.com/paralax/awesome-honeypots}

\bibitem{aws-lambda}
Amazon Web Services, ``AWS Lambda: Deploying Lambda Functions as .zip File Archives,''
\textit{AWS Documentation}. [Online]. Available:
\url{https://docs.aws.amazon.com/lambda/latest/dg/configuration-function-zip.html}

\bibitem{nginx-docs}
NGINX, ``NGINX Documentation,'' \textit{NGINX}. [Online]. Available:
\url{https://nginx.org/en/}

\bibitem{nginx-ec2}
NGINX, ``Create Amazon EC2 Instances for NGINX,'' \textit{NGINX Deployment Guides}.
[Online]. Available:
\url{https://docs.nginx.com/nginx/deployment-guides/amazon-web-services/ec2-instances-for-nginx/}

\bibitem{whatismyip}
WhatIsMyIPAddress, \textit{WhatIsMyIPAddress.com}. [Online]. Available:
\url{https://whatismyipaddress.com}

\bibitem{chatgpt}
OpenAI, ``ChatGPT,'' \textit{OpenAI}. [Online]. Available: \url{https://chatgpt.com/}

\bibitem{eraser}
Eraser, ``AWS Diagram Generator,'' \textit{Eraser.io}. [Online]. Available:
\url{https://www.eraser.io/ai/aws-diagram-generator}

\bibitem{cloudflare-wp}
ResearchGate, ``Enhancing Website Speed and Security With Cloudflare CDN,''
\textit{ResearchGate}, 2024. [Online]. Available:
\url{https://www.researchgate.net/publication/384807865}

\bibitem{eos-ids}
ACM Digital Library, ``An Efficient Optimal Security System for Intrusion Detection
in Cloud Computing Environments Using Hybrid Deep Learning,''
\textit{Advances in Engineering Software}, 2022. [Online]. Available:
\url{https://dl.acm.org/doi/abs/10.1016/j.advengsoft.2022.103236}

\bibitem{rana2024}
K. Rana, ``Intrusion Detection in Cloud Environments Using Hybrid Deep Learning,''
National College of Ireland, 2024. [Online]. Available:
\url{https://norma.ncirl.ie/7728/1/kunalrana.pdf}

\bibitem{fortinet-brute-force}
Fortinet, ``What Is a Brute Force Attack?,'' \textit{Fortinet Cyber Glossary}.
[Online]. Available:
\url{https://www.fortinet.com/resources/cyberglossary/brute-force-attack}

\bibitem{faker}
joke2k, ``Faker Python Package,'' \textit{Faker Documentation}. [Online]. Available:
\url{https://faker.readthedocs.io/en/master/}

\bibitem{locust}
Locust Contributors, ``Locust: An Open Source Load Testing Tool,'' \textit{GitHub}.
[Online]. Available: \url{https://github.com/locustio/locust}

\bibitem{cf-rate-limiting}
Cloudflare, ``What Is Rate Limiting?,'' \textit{Cloudflare Learning Center}.
[Online]. Available:
\url{https://www.cloudflare.com/en-gb/learning/bots/what-is-rate-limiting/}

\bibitem{siteguarding}
SiteGuarding, ``How Hackers Find and Target Vulnerable Websites,''
\textit{SiteGuarding Security Blog}. [Online]. Available:
\url{https://www.siteguarding.com/security-blog/how-hackers-find-and-target-vulnerable-websites-the-complete-technical-breakdown/}

\bibitem{cf-community-blocking}
Cloudflare Community, ``Auto IP Blocking System,'' \textit{Cloudflare Community
Forum}. [Online]. Available:
\url{https://community.cloudflare.com/t/auto-ip-blocking-system/599358}

\bibitem{xvpn}
X-VPN, ``X-VPN Documentation,'' \textit{X-VPN}. [Online]. Available:
\url{https://xvpn.io}

\bibitem{mdn-http-status}
Mozilla Developer Network, ``HTTP Response Status Codes,'' \textit{MDN Web Docs}.
[Online]. Available:
\url{https://developer.mozilla.org/en-US/docs/Web/HTTP/Reference/Status}

\bibitem{burpsuite}
PortSwigger, ``Configuring Burp Intruder Attacks,'' \textit{PortSwigger Web
Security}. [Online]. Available:
\url{https://portswigger.net/burp/documentation/desktop/tools/intruder/configure-attack}

\bibitem{ffuf}
ffuf Contributors, ``ffuf: A Fast Web Fuzzer Written in Go,'' \textit{GitHub}.
[Online]. Available: \url{https://github.com/ffuf/ffuf}

\bibitem{cf-bot-fight}
Cloudflare, ``Get Started with Bot Fight Mode,'' \textit{Cloudflare Bot Management
Documentation}. [Online]. Available:
\url{https://developers.cloudflare.com/bots/get-started/bot-fight-mode/}

\bibitem{cf-verified-bots}
Cloudflare, ``Verified Bots,'' \textit{Cloudflare Radar}. [Online]. Available:
\url{https://radar.cloudflare.com/bots\#verified-bots}

\bibitem{cf-captcha-end}
Cloudflare, ``The End of the Road for Cloudflare CAPTCHAs,'' \textit{Cloudflare
Blog}. [Online]. Available:
\url{https://blog.cloudflare.com/end-cloudflare-captcha/}

\bibitem{ml-clustering-sec}
Uptycs, ``Clustering for Threat Detection,'' \textit{Uptycs Blog}. [Online].
Available:
\url{https://www.uptycs.com/blog/threat-research-report-team/machine-learning-in-cybersecurity}

\bibitem{sklearn-kmeans}
scikit-learn, ``sklearn.cluster.KMeans,'' \textit{scikit-learn Documentation}.
[Online]. Available:
\url{https://scikit-learn.org/stable/modules/generated/sklearn.cluster.KMeans.html}

\bibitem{sklearn-hdbscan}
scikit-learn, ``sklearn.cluster.HDBSCAN,'' \textit{scikit-learn Documentation}.
[Online]. Available:
\url{https://scikit-learn.org/stable/modules/generated/sklearn.cluster.HDBSCAN.html}

\bibitem{sklearn-scaler}
scikit-learn, ``sklearn.preprocessing.StandardScaler,'' \textit{scikit-learn
Documentation}. [Online]. Available:
\url{https://scikit-learn.org/stable/modules/generated/sklearn.preprocessing.StandardScaler.html}

\bibitem{hdbscan-params}
HDBSCAN Contributors, ``Parameter Selection for HDBSCAN*,'' \textit{HDBSCAN
Documentation}. [Online]. Available:
\url{https://hdbscan.readthedocs.io/en/latest/parameter_selection.html}

\bibitem{exploit-db}
Exploit Database, ``Exploit 50702,'' \textit{Exploit-DB}. [Online]. Available:
\url{https://www.exploit-db.com/exploits/50702}

\bibitem{qwen3}
Hugging Face, ``Qwen3-235B-A22B,'' \textit{Hugging Face Model Hub}. [Online].
Available: \url{https://huggingface.co/Qwen/Qwen3-235B-A22B}

\bibitem{hf-inference}
Hugging Face, ``Inference Providers,'' \textit{Hugging Face Documentation}.
[Online]. Available: \url{https://huggingface.co/docs/inference-providers/en/index}

\bibitem{cf-api-lockdown}
Cloudflare, ``Cloudflare API: Firewall Lockdowns,'' \textit{Cloudflare API
Documentation}. [Online]. Available:
\url{https://developers.cloudflare.com/api/resources/firewall/subresources/lockdowns/models/lockdown_url/}

\bibitem{jira-api}
Atlassian, ``JIRA REST API Examples,'' \textit{Atlassian Developer Documentation}.
[Online]. Available:
\url{https://developer.atlassian.com/server/jira/platform/jira-rest-api-examples/}

\bibitem{zero-day}
Wikipedia, ``Zero-Day Vulnerability,'' \textit{Wikipedia}. [Online]. Available:
\url{https://en.wikipedia.org/wiki/Zero-day_vulnerability}

\bibitem{wp-king-addons}
The Hacker News, ``WordPress King Addons Flaw Under Active Attack,''
\textit{The Hacker News}, Dec. 2025. [Online]. Available:
\url{https://thehackernews.com/2025/12/wordpress-king-addons-flaw-under-active.html}

\bibitem{owasp-api8}
OWASP, ``API Security Top 10 API8:2023 Security Misconfiguration,''
\textit{OWASP API Security Project}, 2023. [Online]. Available:
\url{https://owasp.org/API-Security/editions/2023/en/0xa8-security-misconfiguration/}

\bibitem{github-sensitive-data}
GitHub, ``Removing Sensitive Data from a Repository,'' \textit{GitHub Security
Documentation}. [Online]. Available:
\url{https://docs.github.com/en/authentication/keeping-your-account-and-data-secure/removing-sensitive-data-from-a-repository}

\bibitem{nist-800-53}
NIST, ``Security and Privacy Controls for Information Systems and Organizations,''
\textit{NIST SP 800-53 Rev.\ 5}, 2020. [Online]. Available:
\url{https://csrc.nist.gov/publications/detail/sp/800-53/rev-5/final}

\bibitem{mitre-recon}
MITRE, ``Reconnaissance (TA0043),'' \textit{MITRE ATT\&CK Framework}. [Online].
Available: \url{https://attack.mitre.org/tactics/TA0043/}

\bibitem{owasp-waf}
OWASP, ``Web Application Firewall,'' \textit{OWASP Community Controls}. [Online].
Available: \url{https://owasp.org/www-community/controls/Web_Application_Firewall}

\bibitem{concept-drift}
INESC TEC, ``Concept Drift Detection with Clustering via Statistical Change Detection
Methods,'' \textit{INESC TEC Repository}. [Online]. Available:
\url{https://repositorio.inesctec.pt/server/api/core/bitstreams/234d4fc6-af80-4fc2-9d65-2aaddb39508d/content}

\bibitem{cronicle}
J. Huckaby, ``Cronicle: A Multi-Server Task Scheduler,'' \textit{Cronicle}.
[Online]. Available: \url{https://cronicle.net}

\bibitem{azure-frontdoor}
Microsoft, ``Azure Front Door and CDN Documentation,'' \textit{Microsoft Azure
Documentation}. [Online]. Available:
\url{https://learn.microsoft.com/en-us/azure/frontdoor/}

\bibitem{kali-goldeneye}
Kali Linux, ``GoldenEye,'' \textit{Kali Linux Tools}. [Online]. Available:
\url{https://www.kali.org/tools/goldeneye/}

\bibitem{kali-hping3}
Kali Linux, ``hping3,'' \textit{Kali Linux Tools}. [Online]. Available:
\url{https://www.kali.org/tools/hping3/}

\bibitem{tcp-handshake}
O. Sydneyma, ``Understanding the TCP Three-Way Handshake,'' \textit{Medium}.
[Online]. Available:
\url{https://medium.com/@oumasydney2000/understanding-the-tcp-three-way-handshake-a-fundamental-pillar-of-internet-communication-53c9230724c2}

\bibitem{cf-syn-flood}
Cloudflare, ``SYN Flood DDoS Attack,'' \textit{Cloudflare Learning Center}.
[Online]. Available:
\url{https://www.cloudflare.com/en-gb/learning/ddos/syn-flood-ddos-attack/}

\bibitem{cf-ddos}
Cloudflare, ``DDoS Protection,'' \textit{Cloudflare}. [Online]. Available:
\url{https://www.cloudflare.com/en-gb/ddos/}

\bibitem{gartner-cf}
Gartner, ``Cloudflare DDoS Protection Reviews,'' \textit{Gartner Peer Insights}.
[Online]. Available:
\url{https://www.gartner.com/reviews/product/cloudflare-ddos-protection?marketSeoName=ddos-mitigation-solutions}

\bibitem{cf-udp-flood}
Cloudflare, ``UDP Flood DDoS Attack,'' \textit{Cloudflare Learning Center}.
[Online]. Available:
\url{https://www.cloudflare.com/en-gb/learning/ddos/udp-flood-ddos-attack/}

\bibitem{udp-wiki}
Wikipedia, ``UDP Flood Attack,'' \textit{Wikipedia}. [Online]. Available:
\url{https://en.wikipedia.org/wiki/UDP_flood_attack}

\bibitem{kali-slowhttp}
Kali Linux, ``slowhttptest,'' \textit{Kali Linux Tools}. [Online]. Available:
\url{https://www.kali.org/tools/slowhttptest/}

\bibitem{nghttp2-h2load}
nghttp2 Contributors, ``h2load — HTTP/2 Benchmarking Tool,'' \textit{nghttp2
Documentation}. [Online]. Available:
\url{https://nghttp2.org/documentation/h2load-howto.html}

\bibitem{http2-ddos}
L. Ming, Y.-B. Leau, and Y. Xie, ``Distributed Denial of Service Attack in HTTP/2:
Review on Security Issues,'' \textit{IEEE Access}, 2024.
doi: 10.1109/ACCESS.2024.3371013. [Online]. Available:
\url{https://www.researchgate.net/publication/378542947}

\bibitem{bunny-shield-docs}
Bunny.net, ``Bunny Shield Documentation,'' \textit{Bunny.net Documentation}.
[Online]. Available: \url{https://docs.bunny.net/shield}

\bibitem{bunny-waf-docs}
Bunny.net, ``Bunny WAF Documentation,'' \textit{Bunny.net Documentation}.
[Online]. Available: \url{https://docs.bunny.net/shield/waf}

\bibitem{dreamhost-bunny}
DreamHost, ``Understanding the Bunny CDN Metrics Dashboard,'' \textit{DreamHost
Help Center}. [Online]. Available:
\url{https://help.dreamhost.com/hc/en-us/articles/35019587303572-Understanding-the-Bunny-CDN-metrics-dashboard}

\bibitem{cdn-handbook-bunny}
CDN Handbook, ``Bunny CDN,'' \textit{CDN Handbook}. [Online]. Available:
\url{https://www.cdnhandbook.com/providers/bunnycdn/}

\bibitem{european-tech-bunny}
European-Tech, ``Bunny CDN,'' \textit{European-Tech}. [Online]. Available:
\url{https://european-tech.com/service/bunny-cdn/}

\bibitem{logicpin-bunny}
LogicPin, ``BunnyCDN Review 2025,'' \textit{LogicPin}. [Online]. Available:
\url{https://www.logicpin.com/technology/bunnycdn-review-2025-and-how-to-add-it-to-wordpress-developer-friendly-cdn/}

\bibitem{azure-waf-rules}
Microsoft, ``Web Application Firewall DRS and CRS Rule Groups and Rules,''
\textit{Microsoft Azure Documentation}. [Online]. Available:
\url{https://learn.microsoft.com/en-us/azure/web-application-firewall/ag/application-gateway-crs-rulegroups-rules}

\bibitem{stallings2003dos}
W. Stallings, ``Denial of Service Attacks,'' \textit{IEEE Security \& Privacy},
vol.~1, no.~6, pp.~70--73, Nov.-Dec. 2003. doi: 10.1109/MSECP.2003.1253570.

\bibitem{d3fend}
MITRE Corporation, ``D3FEND™ --- A Knowledge Graph of Cybersecurity
Countermeasures,'' \textit{MITRE D3FEND}, 2021. [Online]. Available:
\url{https://d3fend.mitre.org}

\bibitem{aws-gateway}
Amazon Web Services, ``Amazon API Gateway,'' \textit{AWS Documentation}.
[Online]. Available: \url{https://aws.amazon.com/api-gateway/}

\bibitem{aws-cloudwatch}
Amazon Web Services, ``Amazon CloudWatch,'' \textit{AWS Documentation}.
[Online]. Available: \url{https://aws.amazon.com/cloudwatch/}

\bibitem{aws-ec2}
Amazon Web Services, ``Amazon EC2,'' \textit{AWS Documentation}.
[Online]. Available: \url{https://aws.amazon.com/ec2/}

\bibitem{aws-cloudfront}
Amazon Web Services, ``Amazon CloudFront,'' \textit{AWS Documentation}.
[Online]. Available: \url{https://aws.amazon.com/cloudfront/}

\end{thebibliography}

\backmatter

% ============================================================
%  APPENDIX — LOG BOOK
% ============================================================
\appendix

\chapter{Project Log Book}
\label{app:logbook}

\begin{longtable}{p{3cm} p{11cm}}
  \toprule
  \textbf{Date} & \textbf{Tasks Completed / Started} \\
  \midrule
  \endhead
  07 Oct 2025  & Meet with supervisor. \\
  08 Oct 2025  & Research CDNs, DDoS tools, and LOIC. \\
  10 Oct 2025  & Set up event-driven architecture with Lambda, API Gateway, and CloudWatch. \\
  12 Oct 2025  & Build Vue.js frontend for honeypot site. \\
  14 Oct 2025  & Implement SQLite backend for user authentication; add hidden API Gateway link to frontend. \\
  16 Oct 2025  & Spin up EC2 instance; configure security groups; clone Vue project. \\
  20 Oct 2025  & Install Certbot and Nginx on EC2 instance after SSH. \\
  25 Oct 2025  & Analyse initial server logs. \\
  30 Oct 2025  & Set up Cloudflare CDN in front of website. \\
  05 Nov 2025  & Tweak Cloudflare configurations; analyse log effects and attack frequency. \\
  12 Nov 2025  & Start Project Definition Document. \\
  20 Nov 2025  & Write Introduction and Motivation sections. \\
  29 Dec 2025  & Add abbreviations / definitions page. \\
  01 Jan 2026  & Write Research Methodology section. \\
  02 Jan 2026  & Write Software Tools and Architecture sections. \\
  03 Jan 2026  & Further describe honeypot architecture; write up Related Work section. \\
  09 Jan 2026  & Migrate from EC2 to Google Cloud virtual instance. \\
  10 Jan 2026  & Re-configure PM2 and Nginx on GCP VM. \\
  12 Jan 2026  & Fix Cloudflare DNS propagation issue; add definitions to report. \\
  13 Jan 2026  & Re-configure Cloudflare for GCP VM IP. \\
  14 Jan 2026  & Complete Related Work write-up; write rate limit attack script. \\
  15 Jan 2026  & Plan CDN resilience attack tests. \\
  16 Jan 2026  & Run rate limiting tests with different configurations. \\
  17 Jan 2026  & Write up rate limiting test results; continue experimentation. \\
  21 Jan 2026  & Test Cloudflare geographic blocking custom rule. \\
  22 Jan 2026  & Describe geographic blocking results; set up Burp Suite for bot challenge testing. \\
  23 Jan 2026  & Meet supervisor to discuss progress. \\
  24 Jan 2026  & K-Means clustering model training; research ML threat detection. \\
  25 Jan 2026  & Switch to HDBSCAN; write up both clustering algorithms. \\
  28 Jan 2026  & Automate model clustering; add cron job; research LLM interpretability. \\
  01 Feb 2026  & Test LLM with rollup data; retrieve Cloudflare WAF rules via API for LLM context. \\
  02 Feb 2026  & Set up Atlassian account for actionable Jira ticket generation. \\
  04 Feb 2026  & Write up clustering model testing; set up Cronicle dashboard. \\
  05 Feb 2026  & Write up Cronicle setup and usefulness. \\
  06 Feb 2026  & Further research on Cloudflare security features. \\
  07 Feb 2026  & Set up Amazon WAF. \\
  08 Feb 2026  & GoldenEye DDoS attack (Kali Linux). \\
  10 Feb 2026  & Write up DDoS results for Cloudflare. \\
  11 Feb 2026  & DDoS attack using hping3. \\
  12 Feb 2026  & Write up DDoS section; H2Load DDoS test (successful). \\
  13 Feb 2026  & Continue writing DDoS section. \\
  14 Feb 2026  & Complete SYN flood section; write UDP flood section. \\
  18 Feb 2026  & Finish DDoS attacks write-up. \\
  19 Feb 2026  & Set up BunnyCDN. \\
  20 Feb 2026  & Enable Bunny Shield. \\
  27 Feb 2026  & Rate limit test BunnyCDN; DDoS protection test. \\
  01 Mar 2026  & Write up BunnyCDN test results. \\
  02 Mar 2026  & Write up SQL injection results; begin structuring report. \\
  04 Mar 2026  & Label all figures appropriately. \\
  11 Mar 2026  & Organise report sections; switch LLM models. \\
  12 Mar 2026  & Format report sections. \\
  13 Mar 2026  & Finalise report; remove duplicate images. \\
  15 Mar 2026  & Choose open-source LLM inference models. \\
  18 Mar 2026  & Add agentic tooling to LLM. \\
  19 Mar 2026  & Use embedding model to reduce token usage for adding context with references \\
  20 Mar 2026  & Deploy Agentic LLM. \\
  21 Mar 2026  & Read over project report \\
  25 Mar 2026  & Clean up references in project \\
  \bottomrule
\end{longtable}

\end{document}